\chapter{Pochodne i różniczki wyższych rzędów}
\label{ma2:ch:higher-derivatives}

W całym rozdziale $E,F,G$ są rzeczywistymi przestrzeniami
Banacha, a $U\subseteq E$ jest otwarty. Rozważamy odwzorowania
$f:U\to F$, chyba że zaznaczono inaczej. Korzystamy z
analizy funkcjonalnej (AF 15) oraz algebry liniowej (AL 08).
Pochodna rzędu pierwszego jest operatorem liniowym;
kolejne pochodne są naturalnie odwzorowaniami wieloliniowymi.

\section{Przestrzenie odwzorowań wieloliniowych}

\begin{definicja}
\label{ma2:def:multilinear}
Dla $m\ge1$ i przestrzeni unormowanych $E_1,\ldots,E_m,F$
przez $\mathcal L^m(E_1,\ldots,E_m;F)$ oznaczamy przestrzeń
ciągłych odwzorowań $A:E_1\times\cdots\times E_m\to F$
liniowych w każdym argumencie z osobna, z normą
\[
 \|A\|=\sup_{\|h_1\|\le1,\ldots,\|h_m\|\le1}
                 \|A(h_1,\ldots,h_m)\|.
\]
Dla $E_1=\cdots=E_m=E$ piszemy $\mathcal L^m(E;F)$,
a dla $m=1$ także $\mathcal L(E,F)$.
Na skończonym iloczynie przestrzeni unormowanych
stosujemy normę maksimum.
Odwzorowanie takie nazywamy symetrycznym, gdy jego
wartość nie zmienia się przy dowolnej permutacji argumentów.
Podprzestrzeń symetryczną oznaczamy przez $\mathcal L_s^m(E;F)$.
\end{definicja}

Jest to rzeczywista wersja definicji AF 15
(\texttt{def:multilinear15}); algebraiczne podstawy
wieloliniowości i symetrii podano w AL 08
(\texttt{def:lin8-wieloliniowe}).

\begin{stwierdzenie}[Ograniczoność i zupełność]
\label{ma2:prop:multi-banach}
Odwzorowanie $m$-liniowe między przestrzeniami unormowanymi
jest ciągłe wtedy i tylko wtedy, gdy dla pewnej stałej $C\ge0$
\[
 \|A(h_1,\ldots,h_m)\|\le C\prod_{j=1}^m\|h_j\|.
\]
Najmniejszą taką stałą jest $\|A\|$. Jeżeli $F$ jest
przestrzenią Banacha, to $\mathcal L^m(E_1,\ldots,E_m;F)$
jest przestrzenią Banacha.
\end{stwierdzenie}

\emph{Dowód:} AF 15 (\texttt{prop:multi-banach15}).
Ciągłość w zerze daje ograniczenie przez skalowanie
każdego argumentu, a odwrotna implikacja wynika z
wieloliniowego rozwinięcia przyrostu. Zupełność wynika
z punktowej konstrukcji granicy ciągu Cauchy'ego w $F$
i przejścia do granicy w oszacowaniu normy operatorowej.
Sama równoważność ciągłości i ograniczoności nie wymaga
zupełności: dowód źródłowy używa jej dopiero przy
konstrukcji granicy ciągu Cauchy'ego. Minimalność stałej
wynika przez podstawienie argumentów z kul jednostkowych
oraz skalowanie niezerowych argumentów; przy argumencie
zerowym obie strony oszacowania są zerowe.
Własności normy wynikają z nierówności trójkąta
i jednorodności normy w $F$, po wzięciu supremum
na iloczynie kul jednostkowych. Jeśli $\|A\|=0$,
oszacowanie wieloliniowe z $C=0$ daje $A=0$.

\begin{wniosek}
\label{ma2:cor:symmetric-banach}
Jeżeli $F$ jest przestrzenią Banacha, to również
$\mathcal L_s^m(E;F)$ jest przestrzenią Banacha.
\end{wniosek}
\begin{proof}
Jeśli $A_n$ są symetryczne i $A_n\to A$ w normie
wieloliniowej, to dla każdego ustalonego układu argumentów
$A_n(h_1,\ldots,h_m)\to A(h_1,\ldots,h_m)$.
Przechodząc do granicy w równościach dla każdej permutacji,
otrzymujemy symetrię $A$. Podprzestrzeń jest więc domknięta
w przestrzeni Banacha ze stwierdzenia
\ref{ma2:prop:multi-banach}, a zatem zupełna.
\end{proof}

\begin{lemat}[Naturalne utożsamienie izometryczne]
\label{ma2:lem:multi-identification}
Dla $m\ge1$ odwzorowanie
$\mathcal L(E,\mathcal L^m(E;F))\to\mathcal L^{m+1}(E;F)$
dane wzorem
\[
 \widetilde T(h_0,h_1,\ldots,h_m)
       =(Th_0)(h_1,\ldots,h_m)
\]
jest izometrycznym izomorfizmem liniowym.
\end{lemat}

\emph{Dowód:} AF 15 (\texttt{lem:currying15}).
Normy są równe, gdyż supremum po wszystkich argumentach
można obliczyć kolejno po $h_0$ i pozostałych argumentach.
Odwrotność przyporządkowuje $A$ operator
$h_0\mapsto A(h_0,\cdot,\ldots,\cdot)$.
W szczególności dla ustalonych $h_1,\ldots,h_m$
ewaluacja $A\mapsto A(h_1,\ldots,h_m)$ jest liniowa
i ograniczona przez $\prod_j\|h_j\|$.

\section{Iterowane pochodne Frécheta}

\begin{definicja}
\label{ma2:def:ck}
Przyjmujemy $D^0f=f$, $D^1f=Df$. Jeżeli $D^{m-1}f$
istnieje w otoczeniu $p$ i jest różniczkowalne w $p$,
to dla $m\ge2$ definiujemy
\[
 D^mf(p)[h_1,\ldots,h_m]
 =\bigl(D(D^{m-1}f)(p)h_1\bigr)[h_2,\ldots,h_m],
\]
korzystając z lematu \ref{ma2:lem:multi-identification}.
Piszemy $f\in C^m(U,F)$, gdy pochodne wszystkich rzędów
$1,\ldots,m$ istnieją na $U$ i są ciągłe w odpowiednich
normach operatorowych. $C^0(U,F)$ oznacza funkcje ciągłe,
a $C^\infty(U,F)=\bigcap_{m\ge1}C^m(U,F)$.
\end{definicja}

Por. AF 15 (\texttt{def:ck15}). Pierwszy argument
oznacza kierunek ostatniego różniczkowania.
Różniczką rzędu $m$ w $p$ nazywamy odwzorowanie
wieloliniowe $D^mf(p)$, oznaczane także przez
$\mathrm d^mf_p$. Zapis $D^mf(p)[h^m]$ oznacza
wartość na $m$ jednakowych argumentach, nie potęgę operatora.

\begin{twierdzenie}[Symetria drugiej pochodnej -- twierdzenie Schwarza]
\label{ma2:thm:schwarz}
Jeżeli $f\in C^2(U,F)$, to
$D^2f(p)[h,k]=D^2f(p)[k,h]$ dla $p\in U$ i $h,k\in E$.
\end{twierdzenie}

\emph{Dowód:} AF 15 (\texttt{thm:schwarz15}).
Dla dostatecznie małych $s,t$ cały równoległobok
$p+\alpha sh+\beta tk$, $0\le\alpha,\beta\le1$, leży w $U$,
ponieważ jego odległość od $p$ nie przekracza
$|s|\|h\|+|t|\|k\|$. Z dwukrotnego zastosowania
twierdzenia \ref{ma2:thm:integral-increment} różnica
mieszana
\[
 \Delta(s,t)=f(p+sh+tk)-f(p+sh)-f(p+tk)+f(p)
\]
podzielona przez $st\ne0$ ma dwie reprezentacje:
\[
 \int_0^1\int_0^1 D^2f(p+\alpha sh+\beta tk)[k,h]
       \,\mathrm d\beta\,\mathrm d\alpha
\]
oraz
\[
 \int_0^1\int_0^1 D^2f(p+\alpha sh+\beta tk)[h,k]
       \,\mathrm d\alpha\,\mathrm d\beta.
\]
Ciągłość $D^2f$ w $p$ w normie operatorowej
zapewnia jednostajną zbieżność funkcji podcałkowych
do odpowiednich stałych na kwadracie parametrów.
Obie granice muszą być równe. Nie zamieniamy
kolejności całkowania przed udowodnieniem symetrii.

\begin{wniosek}[Symetria wyższych pochodnych]
\label{ma2:cor:higher-symmetric}
Jeżeli $f\in C^m(U,F)$, $m\ge2$, to
$D^mf(p)\in\mathcal L_s^m(E;F)$ dla każdego $p\in U$.
\end{wniosek}

\emph{Dowód:} AF 15 (\texttt{cor:symmetric-higher15}).
Indukcję rozpoczynamy od $m=2$, czyli twierdzenia
\ref{ma2:thm:schwarz}. Dla $m\ge3$ stosujemy je
do $D^{m-2}f\in C^2$ i otrzymujemy symetrię
pierwszych dwóch argumentów.
Różniczkowanie tożsamości symetrii $D^{m-1}f$ po ewaluacji
na ustalonych argumentach daje symetrię ostatnich $m-1$.
Te zamiany sąsiednie generują wszystkie permutacje.
Przestrzeń wartości $\mathcal L^{m-2}(E;F)$ jest zupełna
na mocy stwierdzenia \ref{ma2:prop:multi-banach}.

\begin{twierdzenie}[Symetria drugiej pochodnej w punkcie]
\label{ma2:thm:schwarz-point}
Jeżeli $f$ jest różniczkowalne w otoczeniu punktu
$p\in U$, a $Df$ jest różniczkowalne w $p$, to
$D^2f(p)$ jest symetryczne. Nie zakładamy istnienia
ani ciągłości drugiej pochodnej poza $p$.
\end{twierdzenie}
\begin{proof}
Oznaczmy $A=Df(p)$ i $B=D^2f(p)$. Z definicji
pochodnej w normie operatorowej
\[
 Df(p+v)w=Aw+B[v,w]+R(v)w,\qquad
 \|R(v)\|=o(\|v\|).
\]
Ustalmy $h,k\in E$ i niech $t>0$ będzie tak małe,
aby wszystkie punkty $p+t(sh+rk)$ dla $s,r\in[0,1]$
należały do otoczenia różniczkowalności.
Oznaczmy $M=\|h\|+\|k\|$. Jeśli $M=0$, symetria
wynika z wieloliniowości, więc przyjmijmy $M>0$.
Dla dowolnego $\varepsilon>0$ i dostatecznie małego
$t$ zachodzi $\|R(v)\|\le\varepsilon\|v\|$ dla
$\|v\|\le tM$ (przy $v=0$ mamy $R(0)=0$).
Niech
\[
 q_t(s)=f(p+sth+tk)-f(p+sth)-st^2 B[k,h],
 \qquad 0\le s\le1.
\]
Reguła łańcuchowa daje
\[
 q_t'(s)=t\bigl(R(sth+tk)-R(sth)\bigr)h .
\]
Istotnie, w różnicy pochodnych wyrazy $Ah$
się redukują, a z dwuliniowości pozostaje
$t^2B[k,h]$, który odjęliśmy. W konsekwencji
\[
 \|q_t'(s)\|\le 2\varepsilon t^2M\|h\|.
\]
Funkcja $q_t$ jest ciągła na $[0,1]$ i
różniczkowalna wewnątrz, więc twierdzenie
\ref{ma2:thm:mean-value-inequality} w zastosowaniu
do funkcji parametru $s$ daje
\[
 \left\|\frac{\Delta(t,t)}{t^2}-B[k,h]\right\|
 \le 2\varepsilon M\|h\|,
\]
gdyż $q_t(1)-q_t(0)=\Delta(t,t)-t^2B[k,h]$,
gdzie $\Delta$ jest wcześniejszą różnicą mieszaną.
Z dowolności $\varepsilon$ wynika, że
$\Delta(t,t)/t^2\to B[k,h]$ dla $t\to0^+$.
Powtarzając rozumowanie po zamianie $h,k$,
otrzymujemy granicę $B[h,k]$ tego samego ilorazu.
Jednoznaczność granicy kończy dowód.
\end{proof}

\begin{wniosek}[Punktowa symetria pochodnych dowolnego rzędu]
\label{ma2:cor:pointwise-higher-symmetric}
Niech $m\ge2$. Jeżeli istnieje $D^mf(p)$ w sensie
definicji \ref{ma2:def:ck}, to jest ono symetryczne,
nawet bez założenia $f\in C^m$.
\end{wniosek}
\begin{proof}
Dla $m=2$ stosujemy twierdzenie
\ref{ma2:thm:schwarz-point}. Niech $m\ge3$ i załóżmy
tezę dla rzędu $m-1$. Z definicji $D^{m-1}f$
istnieje w otoczeniu $p$, więc $D^{m-2}f$ jest
tam różniczkowalne, a jego pochodna jest
różniczkowalna w $p$ dzięki istnieniu $D^mf(p)$
i izometrycznemu utożsamieniu z lematu
\ref{ma2:lem:multi-identification}. Punktowe twierdzenie
Schwarza dla $D^{m-2}f$ daje symetrię pierwszych
dwóch argumentów $D^mf(p)$.
Z założenia indukcyjnego $D^{m-1}f(z)$ jest
symetryczne dla każdego $z$ w tym otoczeniu.
Dla dowolnej permutacji $\sigma$ zbioru $\{2,\ldots,m\}$
różniczkujemy w $p$ w kierunku $h_1$ tożsamość
\[
 D^{m-1}f(z)[h_2,\ldots,h_m]
 =D^{m-1}f(z)[h_{\sigma(2)},\ldots,h_{\sigma(m)}].
\]
Ciągłość liniowych ewaluacji zapewnia, że po
różniczkowaniu otrzymujemy tę samą tożsamość
dla $D^mf(p)$ z pierwszym argumentem $h_1$.
Mamy zatem niezmienniczość przy wszystkich
zamianach sąsiednich argumentów, a te generują
wszystkie permutacje.
\end{proof}

Dowody punktowe nie wykorzystują całek ani zupełności,
więc działają również w przestrzeniach unormowanych.
Nie wynika stąd jednak ciągłość $D^mf$ w otoczeniu $p$.

\section{Reguły rachunkowe i przykłady}

\begin{stwierdzenie}[Druga pochodna złożenia]
\label{ma2:prop:second-chain}
Jeżeli $f\in C^2(U,F)$, $g\in C^2(V,G)$, gdzie
$V\subseteq F$ jest otwarty i $f(U)\subseteq V$, to
$g\circ f\in C^2(U,G)$ oraz
\[
 \begin{aligned}
 D^2(g\circ f)(p)[h,k]={}&
 D^2g(f(p))[Df(p)h,Df(p)k]\\
 &+Dg(f(p))(D^2f(p)[h,k]).
 \end{aligned}
\]
\end{stwierdzenie}

\emph{Dowód:} AF 15 (\texttt{prop:second-chain15}).
Stosujemy regułę łańcuchową oraz regułę dla ciągłego
odwzorowania dwuliniowego, tj. składania operatorów,
w twierdzeniach \ref{ma2:thm:frechet-chain}
i \ref{ma2:thm:frechet-bilinear}.

\begin{przyklad}[Druga pochodna funkcji wzdłuż krzywej]
\label{ma2:ex:second-chain-curve}
Niech $I\subseteq\R$ będzie przedziałem otwartym,
$U\subseteq\R^n$ zbiorem otwartym, $g\in C^2(I,U)$
(tj. $g\in C^2(I,\R^n)$ i $g(I)\subseteq U$)
oraz $f\in C^2(U,\R)$. Wtedy
\[
 \begin{aligned}
 (f\circ g)''(t)
 ={}&\sum_{i,j=1}^n
 \partial_i\partial_j f(g(t))g_i'(t)g_j'(t)\\
 &+\sum_{i=1}^n\partial_i f(g(t))g_i''(t).
 \end{aligned}
\]
Ze stwierdzenia \ref{ma2:prop:second-chain},
zastosowanego do złożenia $f\circ g$ i dwóch
kierunków równych $1\in\R$, otrzymujemy
\[
 (f\circ g)''(t)
 =D^2f(g(t))[g'(t),g'(t)]+Df(g(t))g''(t).
\]
Rozwinięcie obu argumentów pierwszego składnika
w bazie standardowej oraz liniowość drugiego dają
wskazane sumy. Równoważnie, różniczkując wzór
z przykładu \ref{ma2:ex:chain-curve}, stosujemy
regułę iloczynu do każdego składnika, a regułę
łańcuchową do funkcji $\partial_j f$.
W notacji macierzowej, z macierzą Hessego
$H_f$ (definicja \ref{ma2:def:hessian}), mamy
\[
 (f\circ g)''(t)=
 g'(t)^{\mathsf T}H_f(g(t))g'(t)
 +\langle\nabla f(g(t)),g''(t)\rangle.
\]

Dla $f(x,y)=x^2y+\sin y$ i $g(t)=(t^2,t)$
mamy $g'(t)=(2t,1)$, $g''(t)=(2,0)$,
$\nabla f(g(t))=(2t^3,t^4+\cos t)$ oraz
\[
 H_f(g(t))=\begin{pmatrix}2t&2t^2\\
                         2t^2&-\sin t\end{pmatrix}.
\]
Pierwszy składnik wzoru wynosi $16t^3-\sin t$,
a drugi $4t^3$. W konsekwencji
$(f\circ g)''(t)=20t^3-\sin t$, zgodnie
z bezpośrednim różniczkowaniem $t^5+\sin t$.

Dla $f(x,y)=x^2+y^2$ i $g(t)=(\cos t,\sin t)$
mamy $H_f=2I_2$, $\nabla f(g(t))=2g(t)$,
$\|g'(t)\|_2=1$ i $g''(t)=-g(t)$.
Pierwszy składnik drugiej pochodnej wynosi $2$,
a drugi $-2$, więc ich suma jest zerowa,
zgodnie z $f\circ g=1$. Sam składnik z macierzą
Hessego nie jest zatem na ogół drugą pochodną
złożenia: trzeba uwzględnić również przyspieszenie
krzywej. Dla krzywej afinicznej $g(t)=p+th$
ten drugi składnik znika.
\end{przyklad}

\begin{stwierdzenie}[Wielomian jednorodny]
\label{ma2:prop:homogeneous-polynomial}
Niech $B\in\mathcal L_s^m(E;F)$, $m\ge1$,
i $P(x)=B(x,\ldots,x)$. Wówczas $P\in C^\infty(E,F)$
i dla $0\le j\le m$
\[
 D^jP(x)[h_1,\ldots,h_j]
 =\frac{m!}{(m-j)!}B(h_1,\ldots,h_j,x,\ldots,x),
\]
gdzie dla $j=0$ lewa strona oznacza $P(x)$.
Pochodne rzędów większych niż $m$ są zerowe.
\end{stwierdzenie}
\begin{proof}
Korzystamy z rozwinięcia wieloliniowego z AF 15
(\texttt{ex:multi-poly15}), tu w przypadku symetrycznym.
Dla $j=0$ wzór jest definicją $P$.
Załóżmy, że wzór zachodzi dla $0\le j<m$.
Dla $j=0$ resztę szacujemy w normie $F$,
a dla $j\ge1$ w normie $\mathcal L^j(E;F)$.
Po podstawieniu $x+u$ w $m-j$ ostatnich argumentach
wyrazy z jednym $u$ dają dokładnie $m-j$ jednakowych
składników. Pozostałe wyrazy z co najmniej dwoma
wystąpieniami $u$ mają, dla $\|h_i\|\le1$ i $\|u\|\le1$,
łączną normę nie większą niż
\[
 \frac{m!}{(m-j)!}\|B\|
 \sum_{\ell=2}^{m-j}\binom{m-j}{\ell}
             \|x\|^{m-j-\ell}\|u\|^\ell ,
\]
z konwencją $a^0=1$. Jest to $O(\|u\|^2)$
jednostajnie względem $h_i$; dla $m-j=1$ suma jest pusta.
Część liniowa w $u$ ma normę ograniczoną przez
$\frac{m!}{(m-j-1)!}\|B\|\|x\|^{m-j-1}\|u\|$.
Daje to pochodną w normie operatorowej i wzór dla $j+1$.
Ciągłość każdej pochodnej wynika z tego samego
rozwinięcia przyrostu (wszystkie wyrazy zawierają $u$).
Dla $j=m$ otrzymujemy stałe odwzorowanie $m!B$,
więc dalsze pochodne są zerowe.
\end{proof}

\begin{przyklad}
W algebrze Banacha niekoniecznie przemiennej
$P(x)=x^3$ ma pochodne
\[
 \begin{aligned}
 DP(x)h&=hx^2+xhx+x^2h,\\
 D^2P(x)[h,k]&=hkx+hxk+khx+xhk+kxh+xkh,\\
 D^3P(x)[h,k,l]&=hkl+hlk+khl+klh+lhk+lkh.
 \end{aligned}
\]
Istotnie, $P(x)=B(x,x,x)$, gdzie
$B$ jest symetryzacją trójliniowego iloczynu
$(a,b,c)\mapsto abc$, tj. średnią sześciu iloczynów
w różnych kolejnościach. Podmultiplikatywność normy
daje $\|B(a,b,c)\|\le\|a\|\|b\|\|c\|$.
Wzory wynikają ze stwierdzenia
\ref{ma2:prop:homogeneous-polynomial}.
\end{przyklad}

\begin{stwierdzenie}[Odtwarzanie odwzorowania symetrycznego z przekątnej]
\label{ma2:prop:polarization}
Niech $B\in\mathcal L_s^m(E;F)$, $m\ge1$, oraz
$P(x)=B(x,\ldots,x)$. Wówczas
\[
 B(h_1,\ldots,h_m)=\frac1{m!}
 \sum_{J\subseteq\{1,\ldots,m\}}(-1)^{m-|J|}
 P\left(\sum_{j\in J}h_j\right).
\]
Suma po pustym zbiorze argumentów oznacza wektor zerowy.
W szczególności $P$ jednoznacznie wyznacza $B$.
\end{stwierdzenie}
\begin{proof}
Jest to wersja wzoru polaryzacyjnego z AL 08
(\texttt{thm:lin8-polaryzacja}) dla wartości w $F$.
Rozwijamy każdy składnik $P$ wieloliniowo.
Ustalony uporządkowany układ $m$ indeksów
ma zbiór występujących indeksów $S$. Jego
współczynnik w sumie przed podzieleniem przez $m!$
jest równy
\[
 \sum_{S\subseteq J\subseteq\{1,\ldots,m\}}(-1)^{m-|J|}
 =(1-1)^{m-|S|}.
\]
Dla $|S|<m$ wynosi on zero na mocy wzoru dwumianowego,
a dla $|S|=m$ wynosi jeden (w sumie jest tylko jeden
składnik). Pozostają więc dokładnie układy, w których
każdy indeks występuje raz. Jest ich $m!$, a symetria
$B$ sprawia, że wszystkie odpowiadają temu samemu
wektorowi $B(h_1,\ldots,h_m)$. Dzielimy przez $m!$.
\end{proof}

Dla $m=2$ wzór przyjmuje postać
\[
 B(h,k)=\tfrac12(P(h+k)-P(h)-P(k)).
\]
Bez symetrii przekątna nie wyznacza odwzorowania:
forma $B((x_1,x_2),(y_1,y_2))=x_1y_2-x_2y_1$
jest niezerowa, chociaż $B(h,h)=0$ dla każdego $h$.
Jest to ważne przy odtwarzaniu wyższych pochodnych
z ich wartości na jednakowych argumentach.

\section{Pochodne mieszane i wielowskaźniki}

Dla $U\subseteq\R^n$ otwartego i $f:U\to F$ zapis
$\partial_i\partial_j f$ oznacza
$\partial_i(\partial_j f)$; pochodną po $j$ liczymy
jako pierwszą. Dla $f\in C^m$ zachodzi
\[
 (\partial_{i_1}\cdots\partial_{i_m} f)(p)
 =D^mf(p)[e_{i_1},\ldots,e_{i_m}],
\]
gdzie $e_i$ to standardowe wektory bazowe.
Istotnie, dla $m=1$ stosujemy twierdzenie
\ref{ma2:thm:frechet-gateaux}. Następnie różniczkujemy
tożsamość dla $m-1$ i używamy reguły łańcuchowej
dla ciągłej liniowej ewaluacji na ustalonych wektorach.
Twierdzenie \ref{ma2:thm:schwarz} i wniosek
\ref{ma2:cor:higher-symmetric} pozwalają zmieniać
kolejność różniczkowania.

\begin{twierdzenie}[Kryterium współrzędnościowe klasy $C^m$]
\label{ma2:thm:ck-partials}
Niech $U\subseteq\R^n$ będzie otwarty, $n\ge1$,
$f:U\to F$ i $m\ge1$. Wówczas $f\in C^m(U,F)$
wtedy i tylko wtedy, gdy wszystkie iterowane
pochodne cząstkowe rzędów od $1$ do $m$ istnieją
na $U$ i są ciągłe.
\end{twierdzenie}
\begin{proof}
Konieczność wynika z poprzedniego wzoru i ciągłości
ewaluacji. Dla $m=1$ dostateczność to wniosek
\ref{ma2:cor:c1-blocks}. Załóżmy tezę dla $m-1$,
$m\ge2$. Wówczas $f\in C^{m-1}$.
Odwzorowanie $A\mapsto(A[e_{i_1},\ldots,e_{i_{m-1}}])$
z $\mathcal L^{m-1}(\R^n;F)$ do $F^{n^{m-1}}$
jest liniowym izomorfizmem topologicznym
(gdzie iloczyn ma normę maksimum).
Istotnie, każda współrzędna jest ograniczona przez
$\|A\|$, a dla dowolnych wektorów $h_j$ rozwinięcie
w bazie standardowej daje
\[
 \|A(h_1,\ldots,h_{m-1})\|
 \le n^{m-1}\max_{i_1,\ldots,i_{m-1}}\|A[e_{i_1},\ldots,e_{i_{m-1}}]\|
             \prod_j\|h_j\|_2.
\]
Odwrotnie, dowolny układ współczynników definiuje
odwzorowanie wieloliniowe przez to rozwinięcie
i spełnia powyższe oszacowanie.
Współrzędne $D^{m-1}f$ mają ciągłe pochodne
cząstkowe rzędu pierwszego na mocy założenia
o pochodnych rzędu $m$. Stosując przypadek $m=1$
do skończonego iloczynu przestrzeni wartości,
otrzymujemy klasę $C^1$ odwzorowania współczynników.
Złożenie z ciągłym liniowym izomorfizmem odwrotnym
daje $D^{m-1}f\in C^1$, czyli $f\in C^m$.
\end{proof}

\begin{definicja}[Wielowskaźnik]
\label{ma2:def:multiindex}
Wielowskaźnikiem nazywamy $\alpha=(\alpha_1,\ldots,\alpha_n)$
o współrzędnych całkowitych nieujemnych. Piszemy
$|\alpha|=\sum_i\alpha_i$, $\alpha!=\prod_i\alpha_i!$,
$h^\alpha=\prod_i h_i^{\alpha_i}$ oraz
$\partial^\alpha f=\partial_1^{\alpha_1}\cdots
\partial_n^{\alpha_n}f$; przy $\alpha_i=0$ odpowiedni
czynnik jest identycznością, a $h_i^0=1$.
\end{definicja}

\begin{wniosek}
\label{ma2:cor:multiindex-diagonal}
Niech $U\subseteq\R^n$ będzie otwarty, $p\in U$.
Dla $f\in C^m(U,F)$, $m\ge1$, i $h\in\R^n$ mamy
\[
 D^mf(p)[h^m]
 =\sum_{|\alpha|=m}\frac{m!}{\alpha!}
                  h^\alpha\,\partial^\alpha f(p).
\]
\end{wniosek}
\begin{proof}
W każdym argumencie rozwijamy $h=\sum_i h_ie_i$.
Wieloliniowość daje sumę po wszystkich uporządkowanych
$m$-elementowych układach indeksów. Symetria pochodnej
sprawia, że wyrazy z tymi samymi licznościami
$\alpha_1,\ldots,\alpha_n$ są jednakowe.
Liczba takich układów wynosi $m!/\alpha!$:
spośród $m!$ permutacji rozróżnialnych pozycji
każdy układ powtarza się $\prod_i\alpha_i!$ razy.
Stosujemy wzór na pochodne mieszane.
\end{proof}

\subsection{Bezpośredni dowód Schwarza dla dwóch zmiennych}

Niech $f:U\subseteq\R^2\to\R$ będzie klasy $C^2$,
$p=(a,b)\in U$. Dla małych niezerowych $h,k$ rozważmy
\[
 \Delta=f(a+h,b+k)-f(a+h,b)-f(a,b+k)+f(a,b).
\]
Prostokąt o tych wierzchołkach zawiera się w $U$.
Twierdzenie Lagrange'a zastosowane do
$x\mapsto f(x,b+k)-f(x,b)$ daje
\[
 \Delta=h\bigl(\partial_1 f(a+\theta h,b+k)
                  -\partial_1 f(a+\theta h,b)\bigr),
 \qquad 0<\theta<1.
\]
Kolejne zastosowanie do $y\mapsto\partial_1 f(a+\theta h,y)$
daje $\Delta/(hk)=\partial_2\partial_1 f
(a+\theta h,b+\eta k)$ dla pewnego $0<\eta<1$.
W odwrotnej kolejności, zaczynając od przyrostu po $y$,
otrzymujemy $\Delta/(hk)=\partial_1\partial_2 f
(a+\theta' h,b+\eta' k)$ z parametrami w $(0,1)$.
Dowód działa również dla ujemnych $h,k$: stosujemy
twierdzenie na przedziale o uporządkowanych końcach,
a znak wynika z różnicy wartości.
Oba punkty pośrednie dążą do $p$ niezależnie od parametrów.
Ciągłość obu pochodnych mieszanych w $p$ daje ich równość.

\begin{przyklad}[Istnienie pochodnych mieszanych nie wystarcza]
\label{ma2:ex:unequal-mixed}
Niech
\[
 f(x,y)=
 \begin{cases}
 xy(x^2-y^2)/(x^2+y^2),&(x,y)\ne(0,0),\\
 0,&(x,y)=(0,0).
 \end{cases}
\]
Funkcja jest gładka poza zerem i ciągła w zerze,
bo $|f(x,y)|\le |xy|\le(x^2+y^2)/2$.
Na osiach obliczamy z definicji
$\partial_2 f(x,0)=x$ oraz $\partial_1 f(0,y)=-y$,
również dla $x=0$ lub $y=0$, ponieważ $f$ znika
na obu osiach. Dokładniej, dla $x\ne0$ mamy
$f(x,t)/t=x(x^2-t^2)/(x^2+t^2)\to x$,
a dla $y\ne0$ analogicznie
$f(t,y)/t=y(t^2-y^2)/(t^2+y^2)\to-y$.
Stąd
$\partial_1\partial_2 f(0,0)=1$,
a $\partial_2\partial_1 f(0,0)=-1$.
Co więcej, $f$ jest klasy $C^1$ na całej płaszczyźnie.
W zerze $Df(0,0)=0$, gdyż $|f(x,y)|\le(x^2+y^2)/2$.
Poza zerem, oznaczając $q=(x^2-y^2)/(x^2+y^2)$,
otrzymujemy z reguły ilorazu
\[
 \partial_1 f=yq+\frac{4x^2y^3}{(x^2+y^2)^2},
 \qquad
 \partial_2 f=xq-\frac{4x^3y^2}{(x^2+y^2)^2}.
\]
Ponieważ $|q|\le1$ oraz
$x^2y^2\le (x^2+y^2)^2/4$, mamy
$|\partial_1 f|\le2|y|$ i $|\partial_2 f|\le2|x|$.
Pochodne cząstkowe dążą więc do zera, a kryterium
\ref{ma2:cor:c1-blocks} daje $f\in C^1$.
Obie pochodne mieszane istnieją w całej płaszczyźnie,
lecz druga pochodna Frécheta w zerze nie istnieje:
gdyby istniała, twierdzenie
\ref{ma2:thm:schwarz-point} wymuszałoby symetrię.
Istotnie, z $\partial_jf(z)=Df(z)e_j$ i ciągłości
liniowej ewaluacji $A\mapsto Ae_j$ wynikałoby
$\partial_i\partial_jf(0,0)=D^2f(0,0)[e_i,e_j]$.
Nie potrzebujemy tu ciągłości drugiej pochodnej.
\end{przyklad}

\section{Macierz Hessego i przekroje}

\begin{definicja}
\label{ma2:def:hessian}
Dla $f\in C^2(U,\R)$, $U\subseteq\R^n$, macierzą Hessego
w $p$ nazywamy $H_f(p)=(\partial_i\partial_j f(p))_{i,j=1}^n$.
\end{definicja}

Z symetrii drugiej pochodnej macierz jest symetryczna.
Rozwijając oba argumenty w bazie standardowej, dostajemy
\[
 D^2f(p)[h,k]=h^{\mathsf T} H_f(p)k.
\]
Dla $h=k$ otrzymujemy formę kwadratową, która będzie
ważna przy badaniu ekstremów. Dla wartości wektorowych
w $\R^q$ mamy $q$ takich macierzy, po jednej dla każdej
funkcji współrzędnej.

\begin{przyklad}
Dla $f(x,y)=x^2y+\sin y$ mamy
\[
 H_f(x,y)=\begin{pmatrix}2y&2x\\2x&-\sin y\end{pmatrix}.
\]
Istotnie, $\partial_1 f=2xy$ i
$\partial_2 f=x^2+\cos y$; różniczkujemy ponownie.
Dla $h=(u,v)$ dostajemy
$D^2f(x,y)[h,h]=2yu^2+4xuv-v^2\sin y$.
Współczynnik $4x$ przy $uv$ wynika z obu równych
wyrazów mieszanych, nie z jednej pochodnej cząstkowej.
\end{przyklad}

\begin{stwierdzenie}[Pochodne przekroju prostoliniowego]
\label{ma2:prop:line-higher}
Niech $m\ge1$, $f\in C^m(U,F)$, $p,h\in E$ i niech
$I=\{t\in\R:p+th\in U\}$.
Funkcja $\gamma:I\to F$, $\gamma(t)=f(p+th)$,
jest klasy $C^m$ i
\[
 \gamma^{(j)}(t)=D^jf(p+th)[h^j],\qquad 1\le j\le m.
\]
\end{stwierdzenie}
\begin{proof}
Zbiór $I$ jest otwarty jako przeciwobraz $U$
przez ciągłe odwzorowanie afiniczne.
Dla $j=1$ teza wynika z reguły łańcuchowej.
Jeśli wzór zachodzi dla $j<m$, różniczkujemy złożenie
$t\mapsto D^jf(p+th)$ z ciągłą liniową ewaluacją
$A\mapsto A(h,\ldots,h)$, której norma nie przekracza
$\|h\|^j$. Reguła łańcuchowa i definicja $D^{j+1}f$
dają wskazany wzór dla $j+1$. Ciągłość wynika z
ciągłości pochodnych $D^jf$ oraz ewaluacji.
\end{proof}

W kolejnym rozdziale zastosujemy stwierdzenie \ref{ma2:prop:line-higher}
do wzoru Taylora z resztą całkową w przestrzeniach Banacha
i do jego zapisu wielowskaźnikowego w $\R^n$.


\section{Przestrzenie funkcji różniczkowalnych}
\label{ma2:sec:differentiable-spaces}

Własności przestrzeni funkcji zależą od wybranej normy.
Niech $I=[-1,1]$, a pochodne na końcach przedziału
oznaczają pochodne jednostronne. Przez afiniczną zmianę
zmiennej wyniki przenoszą się na każdy niedegenerujący
się przedział zwarty. Ośrodkowość oznacza istnienie
przeliczalnego zbioru gęstego.

\begin{lemat}
\label{ma2:lem:continuous-separable}
Przestrzeń $C(I,\R)$ z normą supremum jest ośrodkowa.
\end{lemat}
\begin{proof}
Funkcje odcinkami liniowe na jednorodnych siatkach
o wymiernych wartościach w węzłach tworzą zbiór
przeliczalny. Dla $u\in C(I,\R)$ i $\varepsilon>0$
z jednostajnej ciągłości wybieramy siatkę, na której
oscylacja $u$ na każdym podprzedziale jest mniejsza
niż $\varepsilon/2$. Przybliżamy wartości w węzłach
liczbami wymiernymi z błędem mniejszym niż
$\varepsilon/2$. Ponieważ interpolacja jest kombinacją
wypukłą dwóch wartości w węzłach, błąd na każdym
podprzedziale jest mniejszy niż $\varepsilon$.
\end{proof}

Każda podprzestrzeń przestrzeni metrycznej ośrodkowej
jest ośrodkowa: z przeliczalnej bazy kul o wymiernych
promieniach i środkach w zbiorze gęstym wybieramy po
jednym punkcie z każdego niepustego przecięcia
z podprzestrzenią. Skończony iloczyn przestrzeni
ośrodkowych jest ośrodkowy, gdyż iloczyn ich
przeliczalnych zbiorów gęstych jest przeliczalny i gęsty.

\subsection{Różniczkowalność w ustalonym punkcie}

\begin{definicja}
\label{ma2:def:point-differentiable-space}
Dla $p\in(-1,1)$ definiujemy przestrzeń liniową
\[
 \mathcal D_p(I)=\{f\in C(I,\R):f'(p)\text{ istnieje}\}.
\]
Wymagamy ciągłości na całym $I$, lecz nie
różniczkowalności poza $p$. Rozważamy normy
\[
 \begin{aligned}
 \|f\|_0&=\|f\|_\infty,\\
 \|f\|_{p,1}&=\|f\|_\infty+|f'(p)|,\\
 \|f\|_{p,2}&=|f(p)|+\sup_{x\in I,\ x\ne p}
              \left|\frac{f(x)-f(p)}{x-p}\right|.
 \end{aligned}
\]
\end{definicja}

Pierwsze dwie normy spełniają aksjomaty normy
dzięki liniowości pochodnej i normie supremum.
Dla trzeciej iloraz różnicowy przedłuża się
ciągle do $p$ przez wartość $f'(p)$, więc jest
ograniczony. Liniowość ilorazu i wartości w $p$
daje jednorodność i nierówność trójkąta.
Jeśli norma znika, to $f(p)=0$ oraz $f(x)=f(p)$
dla każdego $x\ne p$, a więc $f=0$.

\begin{twierdzenie}
\label{ma2:thm:point-differentiable-norms}
Przestrzeń $\mathcal D_p(I)$ jest ośrodkowa
we wszystkich trzech normach. W pierwszych dwóch
nie jest zupełna, a jej uzupełnienia to odpowiednio
$C(I,\R)$ i $C(I,\R)\times\R$ z normą sumy.
W trzeciej normie jest przestrzenią Banacha
izometrycznie izomorficzną z $\R\times C(I,\R)$.
\end{twierdzenie}
\begin{proof}
W pierwszej normie mamy izometryczne włożenie do
$C(I,\R)$, w drugiej izometrię $f\mapsto(f,f'(p))$
do $C(I,\R)\times\R$. Ośrodkowość wynika z lematu
\ref{ma2:lem:continuous-separable} i powyższych
własności przestrzeni metrycznych.

Funkcje $f_n(x)=\sqrt{(x-p)^2+1/n^2}$ mają
$f_n'(p)=0$ i zbiegają jednostajnie do $|x-p|$,
gdyż $0\le f_n(x)-|x-p|\le1/n$.
Są zatem ciągiem Cauchy'ego w obu normach,
ale granica jednostajna nie ma pochodnej w $p$.

Aby udowodnić gęstość obrazów, dla dowolnego
$u\in C(I,\R)$ wybieramy coraz drobniejsze
siatki zawierające $p$. Na podprzedziale
$[a,b]$ interpolujemy wartości wzorem
\[
 u(a)+(u(b)-u(a))\rho((x-a)/(b-a)),
 \qquad \rho(s)=3s^2-2s^3.
\]
Mamy $\rho(0)=0$, $\rho(1)=1$,
$\rho'(0)=\rho'(1)=0$ i $0\le\rho\le1$
na $[0,1]$. Interpolanty $u_n$ są więc
klasy $C^1$, mają $u_n'(p)=0$ i zbiegają
jednostajnie do $u$ z jednostajnej ciągłości $u$.
Dla dowolnego $c\in\R$ funkcje
$v_n(x)=u_n(x)+c(x-p)/(1+n^2(x-p)^2)$
mają $v_n'(p)=c$ oraz $v_n\to u$ jednostajnie,
bo $2n|x-p|\le1+n^2(x-p)^2$.
Obrazy są zatem gęste w obu przestrzeniach
docelowych, które są zupełne z twierdzenia
\ref{ma2:thm:continuous} oraz zupełności
skończonego iloczynu.

W trzeciej normie definiujemy ciągłą funkcję
\[
 Q_pf(x)=
 \begin{cases}
 (f(x)-f(p))/(x-p),&x\ne p,\\
 f'(p),&x=p.
 \end{cases}
\]
Odwzorowanie $f\mapsto(f(p),Q_pf)$ jest liniową
izometrią do $\R\times C(I,\R)$ z normą sumy.
Jest surjektywne: parze $(a,q)$ odpowiada
$f(x)=a+(x-p)q(x)$, dla której $f'(p)=q(p)$
z ciągłości $q$. Stąd zupełność i ośrodkowość.
\end{proof}

Normy te nie są równoważne. Trzecia jest zupełna,
a pozostałe nie, co wyklucza równoważność.
Dla $f_n(x)=(x-p)/(1+n^2(x-p)^2)$ mamy
$\|f_n\|_\infty\le1/(2n)$, lecz $f_n'(p)=1$,
więc pierwsza i druga także nie są równoważne.

\subsection{Różniczkowalność na całym przedziale}

\begin{definicja}
\label{ma2:def:everywhere-differentiable-space}
Niech $\mathcal D(I)$ oznacza funkcje rzeczywiste
różniczkowalne w każdym punkcie $I$.
Definiujemy podprzestrzeń
\[
 \mathcal D_b^1(I)=
 \{f\in\mathcal D(I):\sup_{x\in I}|f'(x)|<\infty\}.
\]
Na $\mathcal D(I)$ mamy normę supremum, a na
$\mathcal D_b^1(I)$ również normę
$\|f\|_{1,b}=\|f\|_\infty+\|f'\|_\infty$.
\end{definicja}

Liniowość przestrzeni i własności norm wynikają
z liniowości pochodnej i własności supremum.
Norma supremum jest skończona z ciągłości $f$.
Nie każda funkcja różniczkowalna na $I$ ma
ograniczoną pochodną: dla $f(0)=0$ oraz
$f(x)=x^2\sin(1/x^2)$ przy $x\ne0$ mamy
$f'(0)=0$ i
$f'(x)=2x\sin(1/x^2)-2\cos(1/x^2)/x$.
Przy $x_n=(2\pi n)^{-1/2}$ otrzymujemy
$f'(x_n)=-2/x_n$, więc norma $\|\cdot\|_{1,b}$
nie jest określona na całej $\mathcal D(I)$.

\begin{twierdzenie}
\label{ma2:thm:bounded-differentiable-banach}
Przestrzeń $\mathcal D(I)$ w normie supremum
jest ośrodkowa i niezupełna; jej uzupełnieniem
jest $C(I,\R)$. Przestrzeń $\mathcal D_b^1(I)$
w normie $\|\cdot\|_{1,b}$ jest przestrzenią
Banacha, ale nie jest ośrodkowa.
Dla każdego $p\in I$ norma
$|f(p)|+\|f'\|_\infty$ jest równoważna
$\|\cdot\|_{1,b}$.
\end{twierdzenie}
\begin{proof}
Ośrodkowość, gęstość i niezupełność $\mathcal D(I)$
w normie supremum wynikają z argumentów
interpolacyjnych i kontrprzykładu
$\sqrt{x^2+1/n^2}$ z poprzedniego dowodu.

Niech $(f_n)$ będzie ciągiem Cauchy'ego w
$\|\cdot\|_{1,b}$. Z twierdzeń
\ref{ma2:thm:continuous} i \ref{ma2:thm:bounded}
mamy $f_n\to f\in C(I,\R)$ i $f_n'\to v$
jednostajnie, gdzie $v:I\to\R$ jest ograniczona.
Twierdzenie \ref{ma2:thm:uniform-derivative}
daje $f'=v$ we wnętrzu $I$. Na końcu $p$,
dla dostatecznie dużego $n$ i wszystkich $m\ge n$,
mamy $\|f_m'-f_n'\|_\infty\le\varepsilon$.
Lemat \ref{ma2:lem:vector-increment} dla $f_m-f_n$
na odcinku o końcach $p,x$, po przejściu
$m\to\infty$, daje
\[
 \left|\frac{f(x)-f(p)}{x-p}
 -\frac{f_n(x)-f_n(p)}{x-p}\right|\le\varepsilon .
\]
Drugi iloraz dla ustalonego $n$ dąży do $f_n'(p)$,
a $|f_n'(p)-v(p)|\le\varepsilon$.
Z dowolności $\varepsilon$ wynika $f'(p)=v(p)$
także jednostronnie. Zatem $f\in\mathcal D_b^1(I)$
i $\|f_n-f\|_{1,b}\to0$.
Równoważność norm wynika z
$|f(p)|\le\|f\|_\infty$ oraz
$\|f\|_\infty\le|f(p)|+2\|f'\|_\infty$
na mocy oszacowania przyrostu.

Aby dowieść nieośrodkowości, wybieramy
\[
 \varphi(s)=
 \begin{cases}s(1-s^2)^2,&|s|\le1,\\0,&|s|>1.
 \end{cases}
\]
Funkcja jest klasy $C^1$, gdyż jej wartości
i pochodne w $\pm1$ zgadzają się z zerem.
Mamy $\varphi'(0)=1$. Niech $x_n=2^{-n}$,
$r_n=2^{-2n-3}$ dla $n\ge1$. Przedziały
$[x_n-r_n,x_n+r_n]$ są parami rozłączne:
odległość kolejnych środków to $2^{-n-1}$,
większa od $r_n+r_{n+1}$. Ich jedynym punktem
skupienia jako rodziny jest $0$.
Dla każdego ciągu $a=(a_n)\in\{0,1\}^{\N}$
definiujemy
\[
 f_a(x)=\sum_{n\ge1}a_nr_n
              \varphi((x-x_n)/r_n).
\]
W każdym punkcie co najwyżej jeden składnik
jest niezerowy. Poza zerem suma jest lokalnie
skończona, a $\|f_a'\|_\infty$ poza zerem
nie przekracza $\|\varphi'\|_\infty$.
Dla $x$ z $n$-tego przedziału
$|f_a(x)/x|\le r_n\|\varphi\|_\infty/(x_n-r_n)
\to0$ przy $n\to\infty$, a poza przedziałami
iloraz jest zerowy. Stąd $f_a'(0)=0$,
więc $f_a\in\mathcal D_b^1(I)$.
Jeżeli $a\ne b$, to dla pewnego $n$
$|f_a'(x_n)-f_b'(x_n)|=1$, skąd
$\|f_a-f_b\|_{1,b}\ge1$.
Otrzymaliśmy nieprzeliczalną rodzinę
punktów odległych parami o co najmniej $1$.
Gdyby istniał przeliczalny zbiór gęsty,
każda z rozłącznych kul promienia $1/3$
wokół tych punktów zawierałaby inny jego
element. Jest to sprzeczność.
\end{proof}

\subsection{Przestrzenie funkcji klasy \texorpdfstring{$C^k$}{Ck}}

\begin{twierdzenie}
\label{ma2:thm:ck-banach-separable}
Dla $k\ge0$ przestrzeń $C^k(I,\R)$ funkcji
o ciągłych pochodnych do rzędu $k$ (także
na końcach w sensie jednostronnym) z normą
\[
 \|f\|_{C^k}=\sum_{j=0}^k\|f^{(j)}\|_\infty
\]
jest ośrodkową przestrzenią Banacha.
Dla $k=1$ jest domkniętą właściwą podprzestrzenią
$\mathcal D_b^1(I)$ w jej normie.
\end{twierdzenie}
\begin{proof}
Izometryczne włożenie
$f\mapsto(f,f',\ldots,f^{(k)})$ do
$C(I,\R)^{k+1}$ z normą sumy daje
ośrodkowość. Aby wykazać domkniętość obrazu,
niech wszystkie współrzędne ciągu takich
wektorów zbiegają jednostajnie do ciągłych
funkcji $v_0,\ldots,v_k$.
Dla $j<k$ stosujemy twierdzenie
\ref{ma2:thm:uniform-derivative} do
$f_n^{(j)}$ i $f_n^{(j+1)}$, a na końcach
argument z poprzedniego dowodu.
Otrzymujemy $v_j'=v_{j+1}$ na całym $I$.
Zatem $v_0\in C^k$ i wektor graniczny
należy do obrazu. Zupełność wynika z
zupełności skończonego iloczynu przestrzeni
$C(I,\R)$. Dla $k=1$ jednostajna granica
ciągłych pochodnych jest ciągła, więc
$C^1$ jest domknięta także w $\mathcal D_b^1$.
Jest właściwa: funkcja $u(0)=0$,
$u(x)=x^2\sin(1/x)$ dla $x\ne0$,
ma $u'(0)=0$ i ograniczoną pochodną
$2x\sin(1/x)-\cos(1/x)$ poza zerem,
która nie jest ciągła w zerze
(bierzemy $x_n=1/(2\pi n)$).
\end{proof}

\begin{uwaga}
W przypadku nieskończonego rzędu nie zastępujemy
wszystkich norm $C^k$ ich supremum: może być
nieskończone. Na $C^\infty(I,\R)$ rozważa się
topologię zadaną rodziną norm $\|\cdot\|_{C^k}$,
a na zbiorach otwartych zbieżność jednostajną
wszystkich pochodnych na zbiorach zwartych.
Teorię przestrzeni Frécheta i przestrzeni Schwartza
rozwinięto w AF 17. Normy Höldera opisano
w RRC 16 (\texttt{def:pde16-holder}), a przestrzenie
Sobolewa i ich zupełność w RRC 09
(\texttt{thm:pde9-banach}). Nie utożsamiamy
pochodnych klasycznych z pochodnymi słabymi.
\end{uwaga}


\subsection{Różniczkowalność w punkcie w wielu wymiarach}

Niech $K=[-1,1]^n$, $n\ge1$, $p\in\operatorname{int}K$,
a $F$ będzie rzeczywistą przestrzenią Banacha.
W tym podrozdziale normę w $\R^n$ przyjmujemy
euklidesową. Definiujemy
\[
 \mathcal D_p(K,F)=
 \{f\in C(K,F):Df(p)\text{ istnieje}\}.
\]
Pochodną rozumiemy w sensie Frécheta na wnętrzu $K$.
W odróżnieniu od przypadku jednej zmiennej iloraz
$f(x)/(x-p)$ nie jest tu określony.
Zamiast niego kontrolujemy resztę przybliżenia
liniowego i jej iloraz przez długość przyrostu.

\begin{twierdzenie}[Normy przy różniczkowalności w punkcie]
\label{ma2:thm:point-differentiable-multidim}
Na $\mathcal D_p(K,F)$ określamy normy
\[
 \begin{aligned}
 N_0(f)&=\|f\|_\infty,\\
 N_1(f)&=\|f\|_\infty+\|Df(p)\|,\\
 N_2(f)&=\|f(p)\|+\|Df(p)\|\\
 &\quad+\sup_{x\in K,\ x\ne p}
 \frac{\|f(x)-f(p)-Df(p)(x-p)\|}{\|x-p\|_2}.
 \end{aligned}
\]
W normie $N_2$ jest to przestrzeń Banacha,
izometrycznie izomorficzna z
\[
 F\times\mathcal L(\R^n,F)\times
 \{q\in C(K,F):q(p)=0\},
\]
z normą sumy. Jeżeli $F\ne\{0\}$, to
w normach $N_0,N_1$ przestrzeń nie jest zupełna,
a uzupełnienia to odpowiednio $C(K,F)$
oraz $C(K,F)\times\mathcal L(\R^n,F)$.
We wszystkich trzech normach ośrodkowość
zachodzi wtedy i tylko wtedy, gdy $F$ jest ośrodkowa.
\end{twierdzenie}
\begin{proof}
Dla $a=f(p)$ i $A=Df(p)$ definiujemy
\[
 q_f(x)=
 \begin{cases}
 \dfrac{f(x)-a-A(x-p)}{\|x-p\|_2},&x\ne p,\\
 0,&x=p.
 \end{cases}
\]
Ciągłość w $p$ wynika wprost z definicji pochodnej,
a poza $p$ z ciągłości licznika i niezerowego mianownika.
Odwzorowanie $f\mapsto(a,A,q_f)$ jest liniowe.
Jest bijekcją na wskazany iloczyn: dowolnej trójce
$(a,A,q)$ odpowiada funkcja
\[
 f(x)=a+A(x-p)+\|x-p\|_2q(x).
\]
Istotnie, $f$ jest ciągła, $f(p)=a$ i reszta
po odjęciu $A(x-p)$, podzielona przez $\|x-p\|_2$,
ma normę $\|q(x)\|$, dążącą do zera.
Stąd $Df(p)=A$ oraz $q_f=q$.
Przeniesienie normy sumy daje dokładnie $N_2$,
co dowodzi także aksjomatów normy.
Trzeci czynnik iloczynu jest domknięty w $C(K,F)$,
gdyż ewaluacja $q\mapsto q(p)$ jest ciągła.
Wszystkie czynniki są zupełne na mocy twierdzenia
\ref{ma2:thm:continuous} i stwierdzenia
\ref{ma2:prop:multi-banach}.

Dla $N_0$ i $N_1$ aksjomaty normy wynikają
z normy supremum i liniowości pochodnej.
Odpowiednie izometrie to $f\mapsto f$ oraz
$f\mapsto(f,Df(p))$. Wykażemy gęstość drugiego
obrazu; z tego wynika też gęstość pierwszego.
Niech $u\in C(K,F)$, $A\in\mathcal L(\R^n,F)$
i $\delta>0$ będzie mniejsza od odległości
$p$ od brzegu $K$. Wybieramy ciągłą funkcję
$\theta_\delta:K\to[0,1]$ równą $1$ na
$\|x-p\|_2\le\delta/2$ i równą $0$ dla
$\|x-p\|_2\ge\delta$, na przykład liniową
względem odległości pomiędzy tymi promieniami.
Funkcja
\[
 u_\delta(x)=(1-\theta_\delta(x))u(x)
       +\theta_\delta(x)(u(p)+A(x-p))
\]
jest ciągła i afiniczna blisko $p$, zatem
$Du_\delta(p)=A$. Ponadto
\[
 \|u_\delta-u\|_\infty
 \le\sup_{\|x-p\|_2\le\delta}\|u(x)-u(p)\|
                  +\delta\|A\|\to0.
\]
Obie przestrzenie docelowe izometrii są zupełne.
Gęstość ich obrazów dowodzi twierdzeń o uzupełnieniu.

Jeżeli $v\in F$ ma normę $1$, to
$f_m(x)=\sqrt{(x_1-p_1)^2+m^{-2}}\,v$
jest różniczkowalna w $p$, ma $Df_m(p)=0$
i zbiega jednostajnie do $|x_1-p_1|v$,
która nie jest różniczkowalna w $p$
(ograniczenie do prostej $p+te_1$).
Ciąg $(f_m)$ jest Cauchy'ego w obu normach,
ale nie ma w nich granicy.

Jeśli $F$ jest ośrodkowa, ustalamy przeliczalny
zbiór gęsty $S\subset F$. Na jednorodnych
siatkach kostki $K$ rozważamy funkcje
otrzymane przez interpolację wieloliniową
wartości z $S$ w węzłach. Na każdej małej
kostce interpolacja jest kombinacją wypukłą
wartości w jej $2^n$ wierzchołkach: współczynniki
są iloczynami współczynników liniowej interpolacji
w poszczególnych zmiennych, a ich suma wynosi $1$.
Na wspólnych ścianach formuły są zgodne, gdyż
współczynniki pozostałych wierzchołków znikają.
Rodzina jest przeliczalna, a jednostajna ciągłość
dowolnego $u\in C(K,F)$ pozwala wybrać siatkę
z oscylacją mniejszą od $\varepsilon/2$ na
każdej małej kostce i wartości z $S$
z błędem mniejszym od $\varepsilon/2$.
Błąd interpolacji jest mniejszy od $\varepsilon$.
Stąd $C(K,F)$ jest ośrodkowa.
Przestrzeń $\mathcal L(\R^n,F)$ jest topologicznie
izomorficzna z $F^n$ przez wartości na wektorach
bazowych (oszacowanie normy przez sumę norm
współczynników). Wszystkie wymienione iloczyny
i ich podprzestrzenie są więc ośrodkowe.
Odwrotnie, funkcje stałe $x\mapsto a$, $a\in F$,
tworzą w każdej normie izometryczną kopię $F$,
więc ośrodkowość którejkolwiek z przestrzeni
implikuje ośrodkowość $F$.
\end{proof}

Normę $N_2$ można zastąpić równoważną normą
\[
 N_*(f)=\|f(p)\|+
 \sup_{x\in K,\ x\ne p}
 \frac{\|f(x)-f(p)\|}{\|x-p\|_2}.
\]
Supremum jest skończone z oszacowania reszty
w poprzednim dowodzie i ograniczoności $Df(p)$.
Jest to norma, gdyż iloraz przyrostu jest liniowy
w $f$, a jego znikanie wraz z $f(p)$ daje $f=0$.
Dla dowolnego wektora jednostkowego $h$
punkty $p+th$ należą do $K$ dla małych $t>0$.
Przejście do granicy ilorazu różnicowego pokazuje,
że $\|Df(p)h\|\le N_*(f)-\|f(p)\|$.
Biorąc supremum po $h$, dostajemy
$\|Df(p)\|\le N_*(f)-\|f(p)\|$.
Nierówność trójkąta w obu kierunkach daje
\[
 N_*(f)\le N_2(f)\le3N_*(f).
\]
W szczególności operator $f\mapsto Df(p)$
jest ciągły w obu tych normach, a także w $N_1$.
Nie jest ciągły w $N_0$, gdy $F\ne\{0\}$:
dla $\|v\|=1$ funkcje
$f_m(x)=m^{-1}\sin(m(x_1-p_1))v$
mają normę supremum co najwyżej $1/m$,
ale $\|Df_m(p)\|=1$.
Ewaluacja $f\mapsto f(x)$ jest ciągła we
wszystkich czterech normach, gdyż dla
$N_*$ mamy $\|f(x)\|\le\|f(p)\|
+\|x-p\|_2(N_*(f)-\|f(p)\|)$.

\subsection{Ograniczone pochodne na zbiorach otwartych}

\begin{definicja}
\label{ma2:def:bounded-differentiable-open}
Niech $E,F$ będą rzeczywistymi przestrzeniami
Banacha, $U\subseteq E$ niepusty i otwarty.
Dla $k\ge1$ przez $\mathcal D_b^k(U,F)$
oznaczamy funkcje $f:U\to F$ mające pochodne
Frécheta do rzędu $k$ na $U$, takie że
\[
 \|f\|_{k,b}=\sum_{j=0}^k\sup_{x\in U}\|D^jf(x)\|
 <\infty.
\]
Przez $C_b^k(U,F)$ oznaczamy podprzestrzeń,
w której wszystkie pochodne są ciągłe.
Dla $k=0$ przyjmujemy $C_b^0(U,F)$ jako przestrzeń
ograniczonych funkcji ciągłych.
\end{definicja}

\begin{twierdzenie}
\label{ma2:thm:bounded-differentiable-open}
Przestrzeń $\mathcal D_b^k(U,F)$ z normą
$\|\cdot\|_{k,b}$ jest przestrzenią Banacha.
Przestrzeń $C_b^k(U,F)$ jest jej domkniętą
podprzestrzenią, a więc także przestrzenią Banacha.
Wniosek o zupełności $C_b^0(U,F)$ wynika
z twierdzenia \ref{ma2:thm:bounded} i domkniętości
funkcji ciągłych względem zbieżności jednostajnej.
\end{twierdzenie}
\begin{proof}
Liniowość pochodnych kolejnych rzędów daje
liniowość przestrzeni i, wraz z normą supremum,
aksjomaty normy. W szczególności składnik rzędu
zero zapewnia dodatnią określoność.
Ustalmy ciąg Cauchy'ego $(f_m)$ w $\|\cdot\|_{k,b}$.
Przy konwencji $\mathcal L^0(E;F)=F$,
z twierdzenia \ref{ma2:thm:bounded} i stwierdzenia
\ref{ma2:prop:multi-banach} otrzymujemy granice
$D^jf_m\to A_j$ jednostajnie na $U$,
dla $0\le j\le k$, o wartościach w
$\mathcal L^j(E;F)$. Wykażemy, że $A_j$ są
kolejnymi pochodnymi $A_0$.

Ustalmy $j<k$ i $p\in U$. Wybierzmy kulę
$B(p,r)\subseteq U$. Z lematu
\ref{ma2:lem:multi-identification} traktujemy
$A_{j+1}(p)$ jako operator z $E$ do
$\mathcal L^j(E;F)$ (dla $j=0$ nie potrzeba
utożsamienia). Dla $\varepsilon>0$ wybieramy $m$
tak, że $\sup_U\|D^{j+1}f_l-D^{j+1}f_m\|
\le\varepsilon$ dla wszystkich $l\ge m$
oraz $\sup_U\|D^{j+1}f_m-A_{j+1}\|
\le\varepsilon$. Stosując twierdzenie
\ref{ma2:thm:mean-value-inequality} do
$D^jf_l-D^jf_m$ na odcinku $[p,p+h]$
dla $0<\|h\|<r$ i przechodząc do granicy
$l\to\infty$, dostajemy
\[
 \|(A_j-D^jf_m)(p+h)-(A_j-D^jf_m)(p)\|
 \le\varepsilon\|h\|.
\]
Stąd
\[
 \begin{aligned}
 &\frac{\|A_j(p+h)-A_j(p)-A_{j+1}(p)h\|}{\|h\|}\\
 &\quad\le 2\varepsilon+
 \frac{\|D^jf_m(p+h)-D^jf_m(p)-D^{j+1}f_m(p)h\|}
 {\|h\|}.
 \end{aligned}
\]
W tym wzorze $A_{j+1}(p)h$ i $D^{j+1}f_m(p)h$
oznaczają ustalenie pierwszego argumentu w odwzorowaniu
wieloliniowym. Ostatni iloraz dąży do zera przy
$h\to0$ dla ustalonego $m$. Z dowolności
$\varepsilon$ wynika $DA_j(p)=A_{j+1}(p)$
przy wskazanym utożsamieniu.
Indukcja po $j$ daje $D^jA_0=A_j$.
Wszystkie $A_j$ są ograniczone, zatem
$A_0\in\mathcal D_b^k(U,F)$ i
$\|f_m-A_0\|_{k,b}\to0$.
Jeżeli $f_m\in C_b^k$, to jednostajne granice
ich ciągłych pochodnych są ciągłe, czyli
$A_0\in C_b^k$. Dowodzi to domkniętości.
\end{proof}

Na zbiorze wypukłym $U$ każda funkcja
$f\in\mathcal D_b^1(U,F)$ jest Lipschitzowska
ze stałą $\sup_U\|Df\|$ na mocy twierdzenia
\ref{ma2:thm:mean-value-inequality}.
Odwrotnie, jeśli różniczkowalna funkcja jest
Lipschitzowska ze stałą $L$, to dla każdego
$x\in U$ i wektora jednostkowego $h$ przejście
do granicy ilorazu różnicowego wzdłuż $h$
daje $\|Df(x)h\|\le L$, a więc
$\sup_U\|Df\|\le L$. Ograniczoność wartości $f$
jest osobnym warunkiem, istotnym na zbiorach
nieograniczonych. Zupełność nie wymaga wypukłości
$U$: dowód wykorzystuje jedynie lokalne kule.

\subsection{Ciągłe pochodne aż do brzegu kostki}

\begin{twierdzenie}
\label{ma2:thm:ck-cube-banach-separable}
Niech $K=[-1,1]^n$, $n\ge1$, a $F$ będzie
rzeczywistą przestrzenią Banacha. Przez $C^k(K,F)$
rozumiemy funkcje ciągłe na $K$, klasy $C^k$
na jego wnętrzu, których pochodne do rzędu $k$
przedłużają się ciągle na $K$. Z normą
$\|f\|_{C^k}=\sum_{j=0}^k\sup_K\|D^jf\|$
jest to przestrzeń Banacha, ośrodkowa wtedy
i tylko wtedy, gdy $F$ jest ośrodkowa.
\end{twierdzenie}
\begin{proof}
Przedłużenia są jednoznaczne z gęstości wnętrza.
Zatem odwzorowanie
$f\mapsto(f,Df,\ldots,D^kf)$ jest liniową
izometrią do skończonego iloczynu
$\prod_{j=0}^k C(K,\mathcal L^j(\R^n;F))$
z normą sumy i $\mathcal L^0(\R^n;F)=F$.
Przestrzenie w iloczynie są zupełne.
Jeżeli ciąg wektorów obrazu zbiega w iloczynie
do $(A_0,\ldots,A_k)$, to wszystkie granice
są ciągłe na $K$. Dla $k\ge1$ ograniczenia do
wnętrza $K$ zbiegają w normie z twierdzenia
\ref{ma2:thm:bounded-differentiable-open},
więc $A_j=D^jA_0$ na wnętrzu.
Z ciągłości na $K$ wynika, że są to
wymagane przedłużenia. Dla $k=0$ obraz jest
całym $C(K,F)$. Obraz jest więc domknięty.
Jeżeli $F$ jest ośrodkowa, to
$\mathcal L^j(\R^n;F)$ jest ośrodkowa:
rozwinięcie wieloliniowe w bazie standardowej
daje topologiczny izomorfizm z $F^{n^j}$
(por. dowód twierdzenia \ref{ma2:thm:ck-partials}).
Argument interpolacyjny z dowodu twierdzenia
\ref{ma2:thm:point-differentiable-multidim}
daje ośrodkowość wszystkich czynników.
Ośrodkowość podprzestrzeni metrycznej daje
tezę. Odwrotnie, funkcje stałe tworzą
izometryczną kopię $F$.
\end{proof}

Nie wolno wnioskować, że $C_b^k(\operatorname{int}K,F)$
automatycznie pokrywa się z $C^k(K,F)$.
Na $(-1,1)$ funkcja
$f(t)=(1-t)^2\sin(1/(1-t))$
jest klasy $C_b^1$, ale jej pochodna
$-2(1-t)\sin(1/(1-t))+\cos(1/(1-t))$
nie ma granicy przy $t\to1^-$
(co widać dla argumentów sinusa i cosinusa
$2m\pi$ oraz $(2m+1)\pi$).

\begin{stwierdzenie}
\label{ma2:prop:dbk-cube-nonseparable}
Jeżeli $F\ne\{0\}$, $k\ge1$ i $n\ge1$,
to $\mathcal D_b^k((-1,1)^n,F)$
nie jest ośrodkowa.
\end{stwierdzenie}
\begin{proof}
Korzystamy z rodziny $f_a$ z dowodu twierdzenia
\ref{ma2:thm:bounded-differentiable-banach},
dla której pochodne w punktach $x_m=2^{-m}$
przyjmują wartości $a_m$.
Dla $k=1$ kładziemy $P_a=f_a$.
Dla $k>1$ bierzemy $k-1$ kolejnych funkcji
pierwotnych, za każdym razem z wartością zero
w zerze, korzystając z całki Riemanna funkcji
ciągłej na przedziale zwartym. Otrzymujemy
$P_a^{(k-1)}=f_a$ i $P_a^{(k)}=f_a'$.
Wszystkie pochodne do rzędu $k-1$ są ciągłe,
więc ograniczone na $[-1,1]$; pochodna rzędu
$k$ jest ograniczona z konstrukcji $f_a$.
Dla ustalonego $v\in F$ o normie $1$ definiujemy
$G_a(x)=P_a(x_1)v$. Reguła łańcuchowa
i indukcja dają
\[
 D^jG_a(x)[h_1,\ldots,h_j]=
 P_a^{(j)}(x_1)(h_1)_1\cdots(h_j)_1v
 \quad(1\le j\le k).
\]
W szczególności normy pochodnych są ograniczone,
a ewaluacja $D^kG_a$ na $k$ kopiach $e_1$ w
$(x_m,0,\ldots,0)$ daje $a_mv$.
Różnym ciągom binarnym odpowiadają zatem funkcje
odległe w normie $\|\cdot\|_{k,b}$ o co najmniej $1$.
Argument z rozłącznymi kulami kończy dowód.
\end{proof}


\subsection{Topologie lokalne i przestrzenie Frécheta}

Niech $U\subseteq\R^n$ będzie niepustym zbiorem
otwartym, $n\ge1$, a $F$ rzeczywistą przestrzenią
Banacha. Nie wymagamy ograniczoności funkcji ani
ich pochodnych na całym $U$. Korzystamy z AF 17
(\texttt{def:seminormy-cinf17},
\texttt{lem:exhaustion17},
\texttt{thm:frechet-cinf17}), uogólniając
przypadek skalarny na wartości w $F$.

\begin{definicja}[Topologia lokalna]
\label{ma2:def:ck-local}
Dla $k\in\N_0\cup\{\infty\}$ na $C^k(U,F)$
rozważamy półnormy
\[
 p_{K,j}(f)=\sum_{\ell=0}^j\sup_{x\in K}
                     \|D^\ell f(x)\|,
 \qquad K\subseteq U\text{ zwarty},\quad 0\le j\le k,
\]
gdzie $j$ jest skończoną liczbą całkowitą.
Dla pustego $K$ przyjmujemy $p_{K,j}=0$.
Otoczenia zera wyznaczają skończone przecięcia
zbiorów $\{f:p_{K,j}(f)<\varepsilon\}$,
$\varepsilon>0$. Otoczenia innych punktów
otrzymujemy przez przesunięcie.
\end{definicja}

Skończoność półnorm wynika z ciągłości pochodnych
na zbiorze zwartym, a jednorodność i nierówność
trójkąta z liniowości różniczkowania.
Są to na ogół półnormy, nie normy: funkcja
może znikać w otoczeniu $K$ bez znikania na $U$.
Topologia jest lokalnie wypukła, ponieważ
wskazane otoczenia zera są wypukłe.

\begin{twierdzenie}[Metryka i zupełność]
\label{ma2:thm:ck-local-frechet}
Przestrzeń $C^k(U,F)$ z topologią lokalną
jest przestrzenią Frécheta. Ciąg funkcji zbiega
w niej wtedy i tylko wtedy, gdy wszystkie
pochodne do rzędu $k$ (dla $k=\infty$ każdego
skończonego rzędu) zbiegają jednostajnie
na każdym zwartym podzbiorze $U$.
\end{twierdzenie}
\begin{proof}
Z lematu AF 17 (\texttt{lem:exhaustion17})
wybieramy rosnące wyczerpanie $K_m$ zbiorami
zwartymi, $K_m\subseteq\operatorname{int}K_{m+1}$,
takie że każdy zwarty podzbiór $U$ mieści się
w pewnym $K_m$. Można wziąć
\[
 K_m=\{x\in U:\|x\|_2\le m,\
       \operatorname{dist}(x,\R^n\setminus U)\ge1/m\},
\]
a dla $U=\R^n$ domknięte kule promienia $m$.
Kładziemy $j_m=\min(k,m)$, przy $k=\infty$
rozumiejąc $j_m=m$, oraz $q_m=p_{K_m,j_m}$.
Każda półnorma $p_{K,j}$ jest ograniczona
przez pewną $q_m$, bo można wybrać $m\ge j$
z $K\subseteq K_m$. Rodziny wyznaczają więc
tę samą topologię. Ponadto $q_m$ rosną z $m$
i rozdzielają punkty: jeśli $f(x)\ne0$, to
$x\in K_m$ dla pewnego $m$ i $q_m(f)>0$.

Metrykę zgodną z topologią określa wzór
\[
 d(f,g)=\sum_{m=1}^{\infty}2^{-m}
                        \min(1,q_m(f-g)).
\]
Szereg jest zbieżny przez porównanie z szeregiem
geometrycznym. Symetria i dodatnia określoność
wynikają z własności $q_m$. Nierówność trójkąta
wynika z nierówności
$\min(1,s+t)\le\min(1,s)+\min(1,t)$ dla
$s,t\ge0$. Metryka jest niezmiennicza
względem przesunięć. Warunek $d(f,g)\to0$
implikuje $q_m(f-g)\to0$ dla każdego $m$,
bo rozważamy $m$-ty składnik szeregu.
Odwrotnie, najpierw ograniczamy ogon szeregu
przez $\sum_{m>M}2^{-m}$, a następnie skończoną
sumę przez małość $q_1,\ldots,q_M$.
Ten sam argument dla różnic dwóch elementów
ciągu daje równoważność warunku Cauchy'ego
w metryce i we wszystkich półnormach.
Argumenty te dają także zawierania otoczeń
metrycznych i półnormowych, nie tylko
równoważność zbieżności ciągów.

Niech $(f_i)$ będzie ciągiem Cauchy'ego.
Dla każdego dopuszczalnego $j$ i każdego
niepustego $K_m$ ciąg $D^jf_i|_{K_m}$ jest
jednostajnie Cauchy'ego, ponieważ można zwiększyć
indeks półnormy do $\ell\ge\max(m,j)$.
Z twierdzenia \ref{ma2:thm:continuous}
i stwierdzenia \ref{ma2:prop:multi-banach}
istnieje ciągła granica na $K_m$ o wartościach
w $\mathcal L^j(\R^n;F)$ (dla $j=0$ w $F$).
Jednoznaczność granicy punktowej zapewnia zgodność
na przecięciach. Otrzymujemy funkcje $A_j$ na $U$,
ciągłe, gdyż każdy punkt należy do wnętrza pewnego
$K_m$.

Dla $p\in U$ wybierzmy domkniętą kulę
$\overline B(p,r)\subseteq U$. Mieści się ona
w pewnym $K_m$. Na otwartej kuli pochodne
$f_i$ do dowolnego ustalonego rzędu $j\le k$
zbiegają jednostajnie, więc twierdzenie
\ref{ma2:thm:bounded-differentiable-open}
zastosowane dla rzędu $j\ge1$ daje
$D^\ell A_0=A_\ell$ na tej kuli dla
$\ell\le j$. Dla $k=0$ wystarcza ciągłość $A_0$.
Ponieważ $p$ i $j$ były dowolne,
$A_0\in C^k(U,F)$ oraz $q_m(f_i-A_0)\to0$.
Otrzymaliśmy zupełność metryki i tezę.
\end{proof}

\begin{twierdzenie}[Ośrodkowość topologii lokalnej]
\label{ma2:thm:ck-local-separable}
Dla $0\le k\le\infty$ przestrzeń $C^k(U,F)$
z topologią lokalną jest ośrodkowa wtedy
i tylko wtedy, gdy $F$ jest ośrodkowa.
\end{twierdzenie}
\begin{proof}
Wybierzmy wszystkie niedegenerujące się domknięte
kostki o wymiernych środkach i wymiernych
długościach boków, o bokach równoległych do osi,
zawarte w $U$. Jest to rodzina przeliczalna,
a ich wnętrza pokrywają $U$: dla każdego $x\in U$
można wybrać odpowiednio małą kostkę wokół
wymiernego punktu bliskiego $x$, zawierającą
$x$ we wnętrzu i mieszczącą się w kuli zawartej
w $U$. Oznaczmy tę rodzinę przez $(Q_i)$.

Rozważmy odwzorowanie
\[
 f\longmapsto (D^jf|_{Q_i})_{i,j}
 \quad\text{do}\quad
 \prod_{i,j} C(Q_i,\mathcal L^j(\R^n;F)),
\]
gdzie $j$ przebiega dopuszczalne skończone rzędy.
Jest ono liniowe i injektywne dzięki pokryciu
przez $Q_i$ i składnikowi $j=0$.
Topologia indukowana przez iloczyn jest
topologią lokalną: każde ograniczenie do $Q_i$
jest ciągłe, zaś każdy zwarty $K\subseteq U$
jest pokryty skończenie wieloma wnętrzami $Q_i$,
więc supremum po $K$ jest ograniczone przez
maksimum supremów po tych kostkach.
Otrzymujemy topologiczne włożenie.

Jeżeli $F$ jest ośrodkowa, to każdy czynnik
jest ośrodkową przestrzenią metryczną
z argumentu interpolacyjnego z twierdzenia
\ref{ma2:thm:point-differentiable-multidim}
i izomorfizmu $\mathcal L^j(\R^n;F)\cong F^{n^j}$
dla $j\ge1$. Każdy czynnik ma przeliczalną
bazę otoczeń, a skończone warunki na współrzędne
z tych baz tworzą przeliczalną bazę
przeliczalnego iloczynu. Podprzestrzeń również
ma przeliczalną bazę; wybierając po jednym
punkcie z każdego niepustego zbioru bazowego,
otrzymujemy przeliczalny zbiór gęsty.
Odwrotnie, funkcje stałe tworzą podprzestrzeń
topologicznie izomorficzną z $F$: wszystkie
niezerowe półnormy na funkcji stałej o wartości
$a$ są równe $\|a\|$, a ewaluacja w ustalonym
punkcie $U$ jest odwrotnością tego włożenia.
Z ośrodkowości przestrzeni metrycznej wynika
ośrodkowość tej podprzestrzeni.
\end{proof}

\begin{stwierdzenie}[Brak normowalności]
\label{ma2:prop:ck-local-nonnormable}
Jeżeli $F\ne\{0\}$, to dla każdego $k$
przestrzeń $C^k(U,F)$ z topologią lokalną
nie jest normowalna.
\end{stwierdzenie}
\begin{proof}
Gdyby pewna norma wyznaczała tę topologię,
jej otwarta kula jednostkowa zawierałaby
$\{f:q_M(f)<\varepsilon\}$ dla pewnych
$M,\varepsilon>0$, z rosnącej rodziny półnorm.
Zbiór $K_M$ nie jest całym $U$, gdyż niepusty
otwarty podzbiór $\R^n$ nie jest zwarty
(gdyby był, byłby zarazem otwarty i domknięty,
a $\R^n$ jest spójna i niezwarta).
Wybieramy domkniętą kulę dodatniego promienia
w $U\setminus K_M$ oraz niezerową funkcję
gładką $\psi$ o nośniku w tej kuli.
Istnienie takich funkcji udowodniono w AF 17
(\texttt{lem:cutoffs17}).
Dla $v\in F$, $\|v\|=1$, funkcja $f=\psi v$
jest niezerowa, ale $q_M(f)=0$, więc wszystkie
$t f$, $t\in\R$, należałyby do kuli jednostkowej
domniemanej normy. Przeczy to jej jednorodności.
\end{proof}

Różniczkowanie $D:C^k(U,F)\to
C^{k-1}(U,\mathcal L(\R^n,F))$ dla $k\ge1$
jest liniowe i ciągłe: przy izometrycznych
utożsamieniach kolejnych pochodnych półnorma
$p_{K,j}(Df)$ nie przekracza $p_{K,j+1}(f)$.
Dla $k=\infty$ przyjmujemy $\infty-1=\infty$.
Ewaluacja $D^jf(p)$ jest ciągła, bo jej norma
jest ograniczona przez $p_{\{p\},j}(f)$.
Na $C_b^k$ norma globalna wyznacza topologię
silniejszą od lokalnej, gdyż
$p_{K,j}(f)\le\|f\|_{k,b}$.
Dla funkcji skalarnych ciągłość mnożenia
punktowego w topologii lokalnej wynika
ze wzoru Leibniza: stosując równoważne półnormy
oparte na pochodnych cząstkowych, otrzymujemy
oszacowanie z AF 17
(\texttt{prop:operations-cinf17}).
Równoważność półnorm wynika z rozwinięcia
w bazie standardowej jak w dowodzie twierdzenia
\ref{ma2:thm:ck-partials}.

Na przykład $f_m(x)=x^m$ na $(-1,1)$ zbiega
do zera w lokalnej topologii $C^\infty$:
na $[-r,r]$, $0<r<1$, dla ustalonego $j$ i
$m\ge j$ mamy
$\sup_{|x|\le r}|f_m^{(j)}(x)|\le m^j r^{m-j}\to0$.
Zbieżność ostatniego ciągu wynika z kryterium
ilorazowego. Nie ma jednak zbieżności w normie
supremum na $(-1,1)$, gdyż $\|f_m\|_\infty=1$.

\subsection{Funkcje gładkie na kostce zwartej}

\begin{wniosek}
\label{ma2:cor:cinfty-cube-frechet}
Dla $K=[-1,1]^n$ przestrzeń
$C^\infty(K,F)=\bigcap_{k\ge0}C^k(K,F)$
z topologią wyznaczoną przez normy $C^k$
jest przestrzenią Frécheta, ośrodkową wtedy
i tylko wtedy, gdy $F$ jest ośrodkowa.
Jeżeli $F\ne\{0\}$, nie jest normowalna.
\end{wniosek}
\begin{proof}
Metrykę budujemy jak w twierdzeniu
\ref{ma2:thm:ck-local-frechet}, biorąc
$q_m=\|\cdot\|_{C^m}$. Ciąg Cauchy'ego
ma granicę w każdej z przestrzeni Banacha
$C^m(K,F)$ na mocy twierdzenia
\ref{ma2:thm:ck-cube-banach-separable}.
Granice zgadzają się jako funkcje na $K$,
ponieważ zbieżność w każdej normie implikuje
zbieżność jednostajną. Wspólna granica
należy do $C^\infty(K,F)$ i daje zbieżność
w każdej normie. Ośrodkowość przy ośrodkowym
$F$ wynika z włożenia do przeliczalnego
iloczynu ośrodkowych $C^m(K,F)$, jak w dowodzie
twierdzenia \ref{ma2:thm:ck-local-separable}.
Konieczność wynika z podprzestrzeni funkcji stałych.

Aby wykazać brak normowalności, przypuśćmy,
że kula jednostkowa pewnej zgodnej normy
zawiera $\{f:\|f\|_{C^M}<\varepsilon\}$.
Wybierzmy $v\in F$, $\|v\|=1$, i dla $l\ge1$
zdefiniujmy
\[
 f_l(x)=c l^{-M}\cos(lx_1-(M+1)\pi/2)v,
 \qquad 0<c<\varepsilon/(M+1).
\]
Dla $0\le j\le M$ norma supremum $j$-tej
pochodnej nie przekracza $c l^{j-M}\le c$,
więc $\|f_l\|_{C^M}<\varepsilon$.
Jednak $(M+1)$-sza pochodna w zerze na
kopii wektorów $e_1$ ma normę $cl$.
W konsekwencji $\|f_l\|_{C^{M+1}}\ge cl$.
Ciągłość normy $C^{M+1}$ względem domniemanej
normy oznacza, że jest przez nią ograniczona
z dokładnością do stałego czynnika
(co wynika z ciągłości w zerze i skalowania).
Musiałaby więc być ograniczona na jej
kuli jednostkowej, sprzeczność.
\end{proof}

W dowodach dla otwartego $U$ istotny jest
skończony wymiar dziedziny, który daje
przeliczalne wyczerpanie zbiorami zwartymi.
Dla ogólnej nieskończenie wymiarowej przestrzeni
Banacha nadal można definiować topologię
jednostajnej zbieżności pochodnych na zbiorach
zwartych, ale powyższe dowody metryzowalności,
zupełności i ośrodkowości nie przenoszą się
automatycznie. Natomiast twierdzenie
\ref{ma2:thm:bounded-differentiable-open}
o normie globalnej nie wymaga skończonego wymiaru.


\subsection{Przestrzenie Lipschitza i H\"oldera}

Niech $K=[-1,1]^n$, $n\ge1$, a $F$ będzie
rzeczywistą przestrzenią Banacha. Korzystamy
z warunków Lipschitza i H\"oldera z AM I 16
oraz z normy z RRC 16
(\texttt{def:pde16-holder}). Dla $0<\alpha\le1$
i funkcji $u:K\to F$ kładziemy
\[
 [u]_\alpha=\sup_{x,y\in K,\ x\ne y}
              \frac{\|u(x)-u(y)\|}{\|x-y\|_2^\alpha}.
\]
Wartość może być nieskończona. Dla $\alpha=1$
jest to najmniejsza stała Lipschitza, gdy jest
skończona.

\begin{definicja}
\label{ma2:def:holder-spaces}
Dla $k\ge0$ i $0<\alpha\le1$ definiujemy
\[
 C^{k,\alpha}(K,F)=
 \{f\in C^k(K,F):[D^kf]_\alpha<\infty\},
 \qquad
 \|f\|_{k,\alpha}=\|f\|_{C^k}+[D^kf]_\alpha.
\]
Dla $k=0$ przyjmujemy $D^0f=f$.
Przestrzeń $C^{0,1}(K,F)$ oznaczamy również
przez $\operatorname{Lip}(K,F)$.
\end{definicja}

Przy $k\ge1$ pochodne są rozumiane we wnętrzu
i przedłużone ciągle na $K$, jak w twierdzeniu
\ref{ma2:thm:ck-cube-banach-separable}.
Normy w $\mathcal L^k(\R^n;F)$ są normami
wieloliniowymi. W skończonym wymiarze dziedziny
norma jest równoważna normie opartej na
pochodnych cząstkowych z RRC 16, przez
rozwinięcie w bazie standardowej zarówno
pochodnej, jak i różnicy dwóch pochodnych.

\begin{twierdzenie}[Zupełność]
\label{ma2:thm:holder-banach}
Przestrzeń $C^{k,\alpha}(K,F)$ jest przestrzenią
Banacha. Jeśli $0<\beta<\alpha\le1$, to
włożenie $C^{k,\alpha}(K,F)\to C^{k,\beta}(K,F)$
jest ciągłe.
\end{twierdzenie}
\begin{proof}
Wielkość $[D^kf]_\alpha$ jest półnormą
z liniowości różniczkowania, nierówności
trójkąta i jednorodności. Składnik $C^k$
zapewnia dodatnią określoność.
Niech $(f_m)$ będzie ciągiem Cauchy'ego
w normie $\|\cdot\|_{k,\alpha}$.
Z twierdzenia \ref{ma2:thm:ck-cube-banach-separable}
istnieje $f\in C^k(K,F)$ z $f_m\to f$
w normie $C^k$. Dla $\varepsilon>0$
i dostatecznie dużych $l,m$ mamy
$[D^k(f_l-f_m)]_\alpha\le\varepsilon$.
Dla dowolnych różnych $x,y\in K$
przechodzimy z $l\to\infty$ w nierówności
\[
 \|D^k(f_l-f_m)(x)-D^k(f_l-f_m)(y)\|
 \le\varepsilon\|x-y\|_2^\alpha.
\]
Biorąc supremum po parach, otrzymujemy
$[D^k(f-f_m)]_\alpha\le\varepsilon$.
W szczególności $[D^kf]_\alpha$ jest skończona,
gdyż $[D^kf_m]_\alpha$ jest skończona.
Stąd $f_m\to f$ w pełnej normie.
Ciągłość włożenia dla $\beta<\alpha$
wynika z oszacowania
\[
 [u]_\beta\le(\operatorname{diam}K)^{\alpha-\beta}
 [u]_\alpha,
\]
otrzymanego przez rozdzielenie potęgi mianownika.
\end{proof}

Na $K$ zachodzi $[f]_1=\sup_K\|Df\|$
dla $f\in C^1(K,F)$. Nierówność w jedną
stronę wynika z oszacowania przyrostu
(twierdzenie \ref{ma2:thm:mean-value-inequality})
na odcinkach we wnętrzu, a następnie
z ciągłości $f$ aż do brzegu. Odwrotna
nierówność wynika z ilorazów różnicowych
we wnętrzu oraz ciągłego przedłużenia $Df$.
W szczególności normy Lipschitza i $C^1$
na $C^1(K,F)$ są identyczne.
Samo istnienie skończonej stałej Lipschitza
nie daje jednak różniczkowalności wszędzie,
co pokazuje $f(x)=|x_1|v$ dla $v\ne0$.

\begin{lemat}[Oszacowanie interpolacyjne]
\label{ma2:lem:holder-interpolation}
Jeśli $0<\beta<\alpha\le1$ i $u\in C^{0,\alpha}(K,F)$,
to
\[
 [u]_\beta\le
 (2\|u\|_\infty)^{1-\beta/\alpha}
 [u]_\alpha^{\beta/\alpha}.
\]
\end{lemat}
\begin{proof}
Dla $r=\|x-y\|_2>0$ liczba
$A=\|u(x)-u(y)\|$ spełnia równocześnie
$A\le2\|u\|_\infty$ i $A\le[u]_\alpha r^\alpha$.
Podnosząc pierwsze oszacowanie do potęgi
$1-\beta/\alpha$, drugie do $\beta/\alpha$
i mnożąc, otrzymujemy oszacowanie $A/r^\beta$.
Bierzemy supremum po parach.
Jeżeli któryś z czynników jest zerem,
funkcja jest stała, a teza jest natychmiastowa.
\end{proof}

\begin{twierdzenie}[Nieośrodkowość]
\label{ma2:thm:holder-nonseparable}
Jeżeli $F\ne\{0\}$, to dla wszystkich $k\ge0$
oraz $0<\alpha\le1$ przestrzeń $C^{k,\alpha}(K,F)$
nie jest ośrodkowa.
\end{twierdzenie}
\begin{proof}
Ustalmy $v\in F$ o normie $1$. Wystarczy
skonstruować rodzinę funkcji jednej zmiennej:
$f(x)=u(x_1)v$ ma normy pochodnych i półnormę
H\"oldera równe odpowiednim wielkościom dla $u$
na $[-1,1]$. Nierówność w jedną stronę
wynika z $|x_1-y_1|\le\|x-y\|_2$ i normy
funkcjonału pierwszej współrzędnej równej $1$,
w drugą bierzemy punkty na pierwszej osi
i ewaluujemy pochodną na kopiach $e_1$.

Niech $\phi(s)=(1-s^2)^{k+2}$ dla $|s|\le1$
oraz $\phi(s)=0$ poza tym przedziałem.
Funkcja jest klasy $C^{k+1}$, ponieważ
w $\pm1$ wielomian ma zera rzędu $k+2$.
Jej $k$-ta pochodna nie jest stała na $(-1,1)$.
Wybierzmy $s,t\in(-1,1)$ z $s\ne t$
i $\phi^{(k)}(s)\ne\phi^{(k)}(t)$.
Niech $x_m=2^{-m}$ i $r_m=2^{-2m-3}$.
Przedziały $J_m=[x_m-r_m,x_m+r_m]$
są rozłączne i gromadzą się jedynie w zerze.
Dla $a=(a_m)\in\{0,1\}^{\N}$ definiujemy
\[
 u_a(z)=\sum_{m\ge1}a_m r_m^{k+\alpha}
                         \phi((z-x_m)/r_m).
\]
Poza zerem suma jest lokalnie skończona.
Dla $j\le k$ pochodna na $J_m$ jest
$a_m r_m^{k+\alpha-j}\phi^{(j)}((z-x_m)/r_m)$,
a poza przedziałami i zerem znika.
W zerze wszystkie pochodne do rzędu $k$
istnieją i są zerowe. Istotnie, dla $j=0$
ciągłość wynika z amplitud dążących do zera;
dla $1\le j\le k$ indukcyjnie iloraz
$u_a^{(j-1)}(z)/z$ na $J_m$ ma moduł
co najwyżej
\[
 \frac{r_m^{k+\alpha-j+1}
             \|\phi^{(j-1)}\|_\infty}{x_m-r_m}
 \longrightarrow0.
\]
Ponieważ $r_m/x_m\to0$ i wykładnik
$k+\alpha-j+1>1$, granica jest zerowa.
Poza przedziałami iloraz jest zerowy.
Ciągłość wszystkich tych pochodnych
w zerze wynika z dodatnich wykładników
$k+\alpha-j$. Zatem $u_a\in C^k([-1,1])$.

Funkcja $\phi^{(k)}$ jest ograniczona
i Lipschitzowska na $\R$ (ma ograniczoną
pochodną), a więc ma skończoną globalną
półnormę H\"oldera $M$ dla wykładnika $\alpha$:
dla odległości co najwyżej $1$ stosujemy
stałą Lipschitza, dla większych odległości
dwukrotną normę supremum.
Każda skalowana funkcja
$r_m^\alpha\phi^{(k)}((z-x_m)/r_m)$
ma tę samą półnormę $M$ przez zmianę skali.
Jeśli dwa punkty leżą w jednym $J_m$,
ich różnica wartości $u_a^{(k)}$ spełnia
oszacowanie z $M$. Jeśli leżą w różnych
przedziałach, każdy niezerowy składnik
porównujemy z jego wartością zero na
końcu własnego przedziału między punktami.
Odległości do tych końców nie przekraczają
odległości między punktami. Nierówność
trójkąta daje oszacowanie z $2M$.
Jeżeli tylko jeden punkt leży w przedziale,
stosujemy ten sam argument z jednym końcem;
gdy oba są poza przedziałami, różnica jest zero.
Stąd $[u_a^{(k)}]_\alpha\le2M$.

Dla dwóch różnych ciągów $a,b$ wybieramy $m$
z $a_m\ne b_m$. Iloraz H\"oldera pochodnej
rzędu $k$ różnicy w punktach
$x_m+r_ms$ i $x_m+r_mt$ wynosi
\[
 c=\frac{|\phi^{(k)}(s)-\phi^{(k)}(t)|}
          {|s-t|^\alpha}>0.
\]
Zatem $\|u_a-u_b\|_{k,\alpha}\ge c$
niezależnie od ciągów. Nieprzeliczalna
rodzina rozłącznych kul promienia $c/3$
wyklucza ośrodkowość, jak w dowodzie
twierdzenia \ref{ma2:thm:bounded-differentiable-banach}.
\end{proof}

\subsection{Małe przestrzenie H\"oldera}

Dla $0<\alpha<1$ i ciągłego $u:K\to F$
oznaczamy
\[
 \omega_\alpha(u,\delta)=
 \sup_{\substack{x,y\in K\\0<\|x-y\|_2\le\delta}}
 \frac{\|u(x)-u(y)\|}{\|x-y\|_2^\alpha}.
\]
Jeżeli nie ma takich par, supremum przyjmujemy
równe zero.

\begin{definicja}
\label{ma2:def:little-holder}
Małą przestrzenią H\"oldera nazywamy
\[
 h^{k,\alpha}(K,F)=
 \{f\in C^{k,\alpha}(K,F):
    \lim_{\delta\downarrow0}
        \omega_\alpha(D^kf,\delta)=0\}.
\]
Normę dziedziczymy z $C^{k,\alpha}(K,F)$.
\end{definicja}

\begin{twierdzenie}[Domknięcie funkcji gładkich]
\label{ma2:thm:little-holder-closure}
Dla $0<\alpha<1$ przestrzeń $h^{k,\alpha}(K,F)$
jest domkniętą podprzestrzenią właściwą
$C^{k,\alpha}(K,F)$, jeśli $F\ne\{0\}$,
i jest przestrzenią Banacha.
Jest dokładnie domknięciem $C^\infty(K,F)$
w normie $C^{k,\alpha}$, a także domknięciem
funkcji gładkich w otoczeniu $K$.
Jest ośrodkowa wtedy i tylko wtedy,
gdy $F$ jest ośrodkowa.
\end{twierdzenie}
\begin{proof}
Liniowość i domkniętość wynikają z
nierówności
\[
 \omega_\alpha(D^kf,\delta)
 \le [D^k(f-g)]_\alpha+
              \omega_\alpha(D^kg,\delta).
\]
Dla granicy w normie wybieramy najpierw $g$
z małą pierwszą wielkością, potem małe $\delta$.
Zupełność wynika z twierdzenia
\ref{ma2:thm:holder-banach}.
Dla $f\in C^{k+1}(K,F)$ mamy
$\omega_\alpha(D^kf,\delta)
\le\sup_K\|D^{k+1}f\|\delta^{1-\alpha}$
przez oszacowanie przyrostu na kostce.
Stąd funkcje gładkie należą do małej przestrzeni.
Natomiast dla $v\in F$ o normie $1$ funkcja
$f(x)=(\max(x_1,0))^{k+\alpha}v$
należy do $C^{k,\alpha}$, lecz nie do $h^{k,\alpha}$.
Jej $k$-ta pochodna w kierunku $e_1$ jest
$c(\max(x_1,0))^\alpha v$, gdzie
$c=\prod_{j=1}^k(\alpha+j)>0$
(z pustym iloczynem $1$).
Dla $k\ge1$ pochodne do rzędu $k$ istnieją
także w zerze, co wynika z ilorazów różnicowych
potęg o wykładnikach większych od $1$
przed każdym różniczkowaniem.
Nierówność $|s_+^\alpha-t_+^\alpha|
\le|s-t|^\alpha$ wynika z podaddytywności
potęgi dla $0<\alpha<1$ (AM I 16,
lemat o potęgach H\"oldera).
Dla pary $0,re_1$ iloraz H\"oldera
pochodnej rzędu $k$ wynosi $c$ dla $r>0$.
To dowodzi właściwości podprzestrzeni.

Wykażemy gęstość funkcji gładkich. Niech
$f\in h^{k,\alpha}(K,F)$. Dla $0<\tau<1/2$
kładziemy $b=1-\tau$. Funkcja $f(bx)$ jest
określona na większej otwartej kostce
$(-1/b,1/b)^n$. Wybieramy nieujemną
gładką funkcję $\eta$ o nośniku w kuli
jednostkowej i całce równej $1$.
Istnienie funkcji odcięcia wynika z AF 17
(\texttt{lem:cutoffs17}); normalizujemy przez
dodatnią całkę. Niech
$\eta_\epsilon(y)=\epsilon^{-n}\eta(y/\epsilon)$
dla $0<\epsilon<\tau/4$. Definiujemy
\[
 f_{\tau,\epsilon}(x)=
 \int_{[-\epsilon,\epsilon]^n}
       \eta_\epsilon(y) f(b(x-y))\,\mathrm d\lambda_n(y)
\]
dla $x$ w otoczeniu $K$ tak małym,
aby argumenty $b(x-y)$ pozostawały we
wnętrzu $K$ dla wszystkich $y$ z kostki całkowania.
Jest to całka funkcji ciągłej o wartościach
w przestrzeni Banacha; można ją rozumieć
jako iterowaną całkę Riemanna po współrzędnych
(podstawy całki wektorowej w AM II 14).
Dla $j\le k$ mamy
\[
 D^jf_{\tau,\epsilon}(x)=
 b^j\int\eta_\epsilon(y)D^jf(b(x-y))
                          \,\mathrm d\lambda_n(y).
\]
Różniczkowanie pod całką uzasadnia jednostajna
na zwartych podzbiorach dziedziny ciągłość
odpowiednich pochodnych i oszacowanie
przyrostu z AM II 14, zastosowane do reszty
po odjęciu pochodnej; po podzieleniu przez
normę przyrostu reszta dąży do zera jednostajnie
w zmiennej całkowania. Indukcja daje wzór.
Funkcja $f_{\tau,\epsilon}$ jest nawet gładka:
lokalnie wokół każdego $x$ zmieniamy zmienną
$z=x-y$ i całkujemy $\eta_\epsilon(x-z)f(bz)$
po jednej ustalonej większej kostce, zawartej
w $(-1/b,1/b)^n$ i obejmującej nośnik jądra
dla wszystkich $x$ z tego otoczenia.
Po zmianie zmiennej dodatkowe fragmenty
mają zerowe jądro. Wszystkie pochodne
jądra są ciągłe, a ich ilorazy różnicowe
zbiegają jednostajnie na zwartej kostce
całkowania. Oszacowanie normy całki przez
całkę normy daje różniczkowanie dowolnego rzędu.
Nie wymaga to istnienia wyższych pochodnych $f$.

Gdy $\tau\downarrow0$ i $\epsilon<\tau/4$,
argumenty $b(x-y)$ dążą jednostajnie do $x$
dla $x\in K$ i $y$ z nośnika jądra.
Z jednostajnej ciągłości wszystkich $D^jf$,
$j\le k$, normalizacji jądra oraz $b^j\to1$
wynika $f_{\tau,\epsilon}\to f$ w normie $C^k$.
Ponadto, dla $x,z\in K$ o odległości
co najwyżej $\delta$, wzór całkowy daje
\[
 \omega_\alpha(D^kf_{\tau,\epsilon},\delta)
 \le b^{k+\alpha}\omega_\alpha(D^kf,b\delta)
 \le\omega_\alpha(D^kf,\delta).
\]
Wynika to z nieujemności jądra, jego
całki równej $1$ i faktu, że odległość
argumentów po przesunięciu i skalowaniu
wynosi $b\|x-z\|_2$.
Dla różnicy pochodnych rzędu $k$ ilorazy
dla odległości co najwyżej $\delta$ są
ograniczone przez $2\omega_\alpha(D^kf,\delta)$,
a dla większych odległości przez
\[
 2\|D^k(f_{\tau,\epsilon}-f)\|_\infty/\delta^\alpha.
\]
Najpierw wybieramy małe $\delta$, potem
małe $\tau$ i $\epsilon$. Dostajemy zbieżność
w normie $C^{k,\alpha}$, co dowodzi
charakterystyki przez domknięcie.

Jeśli $F$ jest ośrodkowa, to $C^{k+1}(K,F)$
jest ośrodkowa z twierdzenia
\ref{ma2:thm:ck-cube-banach-separable}.
Włożenie do $C^{k,\alpha}$ jest ciągłe,
bo $[D^kf]_\alpha\le
(\operatorname{diam}K)^{1-\alpha}\sup_K\|D^{k+1}f\|$.
Obraz przeliczalnego zbioru gęstego jest
gęsty w obrazie tej przestrzeni, a jego
domknięcie zawiera małą przestrzeń przez
wykazaną aproksymację funkcjami gładkimi.
Jednocześnie całe $C^{k+1}$ należy do
małej przestrzeni, która jest domknięta.
Stąd otrzymujemy przeliczalny zbiór gęsty.
Odwrotnie, funkcje stałe tworzą izometryczną
kopię $F$ w małej przestrzeni.
\end{proof}

Dla $\alpha=1$ nie przenosimy powyższej
charakterystyki: warunek
$\omega_1(u,\delta)\to0$ na wypukłej kostce
wymusza stałość $u$. Istotnie, dzieląc
odcinek od $x$ do $y$ na części długości
co najwyżej $\delta$, otrzymujemy
$\|u(x)-u(y)\|\le\omega_1(u,\delta)\|x-y\|_2$.
Przechodząc do granicy, dostajemy $u(x)=u(y)$.
Tymczasem w domknięciu funkcji gładkich
w normie Lipschitza znajdują się już
niestałe funkcje afiniczne. Dlatego w tym
podrozdziale małe przestrzenie H\"oldera
rozważamy tylko dla $0<\alpha<1$.

\subsection{Norma globalna na otwartej kostce}

\begin{stwierdzenie}
\label{ma2:prop:cbk-open-nonseparable}
Dla $F\ne\{0\}$ oraz każdego skończonego $k\ge0$
przestrzeń $C_b^k((-1,1)^n,F)$ nie jest ośrodkowa.
\end{stwierdzenie}
\begin{proof}
Wykorzystujemy funkcję $\phi$ klasy $C^{k+1}$
z dowodu twierdzenia \ref{ma2:thm:holder-nonseparable}
oraz $v\in F$ o normie $1$. Przenosimy
rozłączne przedziały do punktów
$x_m=1-2^{-m}$ z promieniami $r_m=2^{-2m-3}$.
Ich jedyny punkt skupienia leży poza otwartym
przedziałem. Dla $a\in\{0,1\}^{\N}$ funkcja
\[
 f_a(x)=\sum_{m\ge1}a_m r_m^k
                  \phi((x_1-x_m)/r_m)v
\]
jest lokalnie skończoną sumą funkcji $C^k$
na otwartej kostce. Pochodna rzędu $j\le k$
ma normę co najwyżej
$\|\phi^{(j)}\|_\infty$, gdyż w każdym
punkcie co najwyżej jeden składnik
jest niezerowy i $r_m^{k-j}\le1$.
Dla $s\in(-1,1)$ z $\phi^{(k)}(s)\ne0$
ewaluacja pochodnej rzędu $k$ w punkcie
$(x_m+r_ms,0,\ldots,0)$ na kopiach $e_1$
daje $a_m\phi^{(k)}(s)v$ (dla $k=0$
jest to wartość funkcji).
Różne ciągi binarne dają więc funkcje
odległe w normie globalnej o co najmniej
$|\phi^{(k)}(s)|>0$. Jak poprzednio
wyklucza to ośrodkowość.
\end{proof}
